% concatenated from paper_exSDIRK.tex
\documentclass[10pt,a4paper]{article}
\usepackage{amsmath}                    % From AMS/CTAN distribution
\usepackage{amssymb}                    % From AMS/CTAN distribution
\usepackage{amsthm}                     % From AMS/CTAN distribution
\usepackage{amsfonts}                   % From AMS/CTAN distrib
\usepackage{subcaption}
\usepackage[usenames, dvipsnames]{color}
%\usepackage[dvips]{graphicx}
\usepackage{graphics,graphicx}
%\usepackage{showlabels}
\usepackage{enumerate}
%\usepackage{epsfig}
%\usepackage{GGGraphics}
\usepackage{multirow}
\usepackage{algorithm}
\usepackage{algpseudocode}
\usepackage{dsfont}
\usepackage{float}
\usepackage{mathtools}
\usepackage{bbm,amsbsy}
\usepackage[hidelinks]{hyperref}
\usepackage{lipsum}
%=============================================================================================
%=============================================================================================


\allowdisplaybreaks


\newtheorem{theorem}{Theorem}[section]
\newtheorem{definition}[theorem]{Definition}
\newtheorem{proposition}[theorem]{Proposition}
\newtheorem{lemma}[theorem]{Lemma}
\newtheorem{remark}[theorem]{Remark}
\newtheorem{corollary}[theorem]{Corollary}
%\newtheorem{algo}[theorem]{Algorithm}
\newtheorem{ex}[theorem]{Example}

%%% reduce margins
%\textheight 22 cm \topmargin -4.4 mm 
%\textwidth 16 cm \oddsidemargin -2.6 mm 
%\evensidemargin =\oddsidemargin

\addtolength{\hoffset}{-1.4cm} 
\addtolength{\textwidth}{2.8cm}
\addtolength{\headsep}{0.5cm}
\addtolength{\voffset}{-0.5cm}%

\def\ts{\thinspace}
\DeclareMathOperator{\arcsinh}{arcsinh}
\DeclareMathOperator{\diag}{diag}
\DeclareMathOperator{\sinc}{sinc}
\renewcommand{\epsilon}{\varepsilon}
\newcommand{\eps}{\varepsilon}
\newcommand{\parenthnewln}[1]{\right.\\#1&\left.{}}
\newcommand{\R}{\mathbb{R}}
\newcommand{\bigo}{\mathcal{O}}
\newcommand{\CC}{\mathcal{C}}
\newcommand{\cN}{\mathcal{N}}
\newcommand{\IP}{\mathbb{P}}
% space is R^d or torus T^d
\def \IR{\mathbb R}
\newcommand{\EE}{\IR^d}
\def \IN{\mathbb N}
\def \IC{\mathbb C}
\def \IQ{\mathbb Q}
\def \ZZ{\mathbb Z}
\def \IE{\mathbb E}
\def \IF{\mathbb F}
\def \IT{\mathbb T} 
\def \eps{\epsilon}
\newcommand{\moda}[1]{\textcolor{blue}{#1}\index {#1}}
\newcommand{\modg}[1]{\textcolor{red}{#1}\index {#1}}
\newcommand{\modi}[1]{\textcolor{ForestGreen}{#1}\index {#1}}
\DeclareMathOperator{\sign}{sign}

\def \H{\mathcal H}
\def \s{\rm s}

\def \calA{\mathcal{A}}
\def \calB{\mathcal{B}}
\def \calC{\mathcal{C}}
\def \calJ{\mathcal{J}}
\newcommand{\calN}{\mathcal{N}}
\newcommand{\calL}{\mathcal{L}}
\newcommand{\calS}{\mathcal{S}}
\newcommand{\D}{\mathcal{D}}
\renewcommand{\R}{\mathbb{R}}
%new macros
\newcommand{\dif}[2]{\frac{\partial #1}{\partial #2}}
\newcommand{\diff}[2]{\frac{\partial^2 #1}{\partial #2^2}}
\newcommand{\norm}[1]{\left\lVert #1 \right\rVert}


\hyphenation{Lip-schitz-ness}

\usepackage[symbol]{footmisc}

\renewcommand{\thefootnote}{\fnsymbol{footnote}}


%\input{ex_shared}

\title{
	Explicit stabilized implementation of singly diagonally implicit Runge-Kutta methods
}

\author{ 
	Ibrahim Almuslimani\footnotemark[1], Gilles Vilmart\footnotemark[2], and Konstantinos Zygalakis\footnotemark[3]
}

\begin{document}

\maketitle
\footnotetext[1]{\'Ecole Polytechnique F\'ed\'erale de Lausanne (EPFL), Swiss Plasma Center (SPC), CH-1015 Lausanne, Switzerland. Ibrahim.Almuslimani@epfl.ch}
\footnotetext[2]{Section de Math\'ematiques, Universit\'e de Gen\`eve, CP 64, 1211 Gen\`eve 4, Switzerland. Gilles.Vilmart@unige.ch}
\footnotetext[3]{School of Mathematics and the Maxwell Institute for Mathematical Sciences, University of Edinburgh, Edinburgh, EH9 3FD, UK. K.Zygalakis@ed.ac.uk}
% REQUIRED
\begin{abstract}
Implicit methods are a natural approach for the integration of stiff differential equations, to avoid time-step restrictions faced by standard explicit integrators. Explicit stabilised integrators are an alternative to implicit methods, which can be particularly efficient in high-dimensional applications with diffusive terms. Towards the best of both worlds, we introduce a new explicit stabilised implementation of a class of diagonally implicit Runge-Kutta methods. This allows us to implement high-order singly-diagonally-implicit Runge-Kutta methods for advection-diffusion-reaction PDEs with a provable computational cost analogous to that of standard explicit stabilised methods. The main ingredient is to recast the implicit Runge–Kutta update as the steady state of a modified auxiliary system, which is then computed using a partitioned Runge–Kutta–Chebyshev method inspired by optimisation techniques.

\smallskip
%
\noindent
{\it Keywords:\,}
diagonally implicit Runge-Kutta methods, explicit stabilised methods, partitioned Runge-Kutta methods, stiff problems, advection-diffusion-reaction problems.
\smallskip

\noindent
{\it AMS subject classification (2010):\,}
65L20, 65M12, 65M20
\end{abstract}

 \section{Introduction}

In this paper, we consider systems of ordinary differential equations (ODEs) representing space discretisations of partial differential equations (PDEs) of the form
\begin{subequations}\label{eq:main}
\begin{eqnarray}
    \dot{y} &=& F(y), \\
    F(y) &:=& F_{D}(y)+F_{A}(y)+F_{R}(y).
\end{eqnarray}
\end{subequations}
where $F_{D}(y),F_{A}(y),F_{R}(y) \in \mathbb{R}^{d}$ represent diffusion, advection and reaction terms respectively. 
Typically \eqref{eq:main} arises from the spatial discretisation of an advection-diffusion-reaction problem in $N$ spatial dimensions, and assuming a spatial grid of size $\Delta x$, the dimension of the ODE grows like $d =\mathcal{O}((\Delta x)^{-N})$. 

The numerical time integration of \eqref{eq:main} is very challenging since this is a \emph{stiff} ODE. {Following \cite{HaW96}}, we call a differential equation \emph{stiff} when
%the time step one would choose for a prescribed accuracy it is much greater than the time step allowed by 
{a standard explicit numerical integrator such as a standard explicit Runge-Kutta method faces a severe time-step restriction}.  In particular, in the case of \eqref{eq:main} the eigenvalues of the Jacobian of $F_{D}$ are typically distributed along the negative real axis in an interval that grows as $[-\mathcal{O}(\Delta x^{-2}),0]$ (for a symmetric diffusion operator), which implies that the time step $\Delta t$ that can be used by the explicit Euler method {suffers from the CFL} $\Delta t \leq C \Delta x^{2}$. Furthermore, the eigenvalues of the Jacobian of $F_{A}$ are typically distributed along the {imaginary} axis in an interval that grows as  $[- i\mathcal{O}(\Delta x^{-1}), i \mathcal{O}(\Delta x^{-1})]$, while the eigenvalues of the Jacobian of the reaction term $F_{R}$ even though are usually not related to $\Delta x$ can have a ratio $\max_{j}|\mathcal{R} \lambda_{j}|/\min_{j}|\mathcal{R} \lambda_{j}|$ that can be very large or vary over several orders of magnitude {when modeling multiscale reaction terms}. 


        
%Differential equations hold a central place in the field of mathematical modeling \cite{} as they are the natural mathematical language to describe various physical phenomena.  Furthermore, in the last 50 years, there has been a tremendous amount of work developed around solving numerically differential equations. In particular, today we have a high number of specialized numerical methods that are able to deal with the individual properties of the differential equation one would like to solve. Examples include explicit and implicit Runge-Kutta methods that sometimes come with special features such as symplecticity and/or volume preservation, as well as multistep methods \cite{HNW93,HaW96,HLW02}.

%One of the most important features of interest for a numerical method is the one of \emph{accuracy}. As the name implies \emph{accuracy} relates to how accurately we can approximate the true solution of our differential equation. This problem is normally interconnected with the ones of \emph{numerical stability} and \emph{stiffness}. 

There are two different approaches to dealing with the issue of stiffness within the class of Runge-Kutta methods. The first relates to the design of implicit methods, with a prime example being the implicit Euler method. These methods have the %property of being able to take arbitrary large time steps 
{advantage of achieving stability unconditionally with respect to the stepsize,} 
which comes with the cost of solving a system of nonlinear equations.  When the dimension of the corresponding differential equation is low, the nonlinear system can be solved {efficiently using Newton's method \cite{NW06}}. As the dimension of the differential equation increases, diagonally implicit Runge-Kutta methods are attractive as they lead to a triangular system which is easier to solve using lower-dimensional Jacobian matrices using an appropriate partitioning of the coefficient matrix as proposed in \cite{VSC92}. Nevertheless, when one deals with differential equations like \eqref{eq:main}, dealing with the corresponding nonlinear equations (even in the triangular form) can be cumbersome as it requires specialised linear algebra techniques, {where the use of preconditioners raises issues with both memory access and storage on parallel supercomputers.} 


An alternative approach to using implicit methods for solving \emph{stiff} differential equations is to use explicit stabilised methods \cite{Abdrev}. In their basic form, these are first-order explicit methods that can take larger time steps than the standard Euler method (for an easily computable numerical cost). This property together with the simplicity of implementation has made explicit stabilised methods the method of choice for diffusion-dominated problems. In particular, the second-order Runge Kutta Chebyshev (RKC) method was very successful for large dimensional parabolic problems, while methods with nearly optimal extended stability domain sizes were proposed as the second-order ROCK2 method \cite{AbM01} and the fourth-order ROCK4 method \cite{Abd02}. In the spirit of implicit-explicit methods {(IMEX) \cite{ARW95,ARW97}}, the versatility of explicit stabilised methods allowed many extensions to partitioned Runge-Kutta methods for addressing problems with additional advection and possibly stiff reaction terms. Examples include {Implicit Runge Kutta Chebyshev (IRKC)  \cite{SSV06}} to deal with stiff reaction terms and Partitioned Runge Kutta Chebyshev (PRKC) \cite{Zb11} to include advection terms. In addition, the Partitioned Implicit Orthogonal Runge Kutta Chebyshev method  (PIROCK) \cite{AV13} combined these two approaches to propose a second-order method for solving advection-diffusion reaction equations.  
In particular, PIROCK has recently emerged as an efficient alternative to classical splitting approaches for solar atmosphere simulations in astrophysics \cite{WVMP24} and for large-scale gyrokinetic simulations in plasma physics \cite{AUB26,AOUB26}.
Another recent use of explicit stabilised methods is for solving optimisation problems. More precisely, in  \cite{EVVZ21} by exploiting the good stability properties of explicit stabilised methods a new method called the \emph{Runge–Kutta Chebyshev} descent (RKCD) was proposed and showed to perform similarly to the state of the art for strongly convex optimisation problems.

A major limitation of methods such as IRKC, PRKC, and PIROCK that use an explicit stabilised method to deal with the $F_{D}$ term in \eqref{eq:main}, is that they {are} at most second-order accurate in $\Delta t$. {Indeed, while the construction of the PIROCK method which is based on a splitting into three terms needed to solve 12 order conditions, similar construction of higher order methods would be cumbersome due to the dramatic increase of order conditions for partitioned Runge-Kutta methods. }
On the other hand, implicit methods can achieve 
much higher order of accuracy in $\Delta t$, 
but as discussed above the implementation as 
well as the storage requirements of {the Newton 
method and variants} in high dimensions become an issue.  
In this paper, inspired by the success of RKCD 
as an optimisation method, we revisit diagonally implicit 
methods and their implementation. In 
particular, instead of solving the nonlinear 
system with variants of Newton's method, we 
introduce a new variant of the RKCD method 
that only {outputs} first-order information.  This method is simple to implement due to its explicit nature which circumvents the 
need for specialised linear algebra tools.
%, which become delicate for large-dimensional problems (preconditioners, memory requirements, parallel implementations). 
Furthermore,  it leads to rigorous estimates of the computational complexity of the proposed method. In particular, we can show that
\begin{enumerate}
\item In the pure linear diffusion case ($F_{A}=F_{R}=0$ in \eqref{eq:main}), { we prove that one time step} of an {SDIRK} method\footnote{ Here $\Delta t$ is the time step used by the {SDIRK} method.} (of arbitrary order) by an iterative method requires  $\mathcal{O}(\sqrt{\Delta t}/\Delta x)$ evaluations of $F_{D}$ avoiding the standard CFL condition of standard explicit Runge Kutta methods. This essentially implies that we are able to implement a higher order diagonally implicit Runge-Kutta method at the cost of an explicit stabilised method. 
\item {For a class of linear advection-diffusion PDE ($F_{R}=0$), we prove that in spite of relatively large Peclet numbers the proposed partitioning allows to implement one step of an SDIRK method with a similar complexity to the one of the pure diffusion case discussed at the previous point}.
\item {For various nonlinear diffusion-advection-reaction problems, numerical experiments illustrate the robustness of the approach where we show that the explicit stabilised iterations converge in the case of an order 4 SDIRK method with error control, a state-of-the-art Runge-Kutta method proposed in \cite[Table IV.6.5, p. 100]{HaW96} and considered in this paper to illustrate the new approach.}
\end{enumerate}






The paper is organized as follows. In Section \ref{sec:pre} we revisit the ideas of numerical stability of Runge-Kutta methods, while also discussing the basics of explicit stabilised methods. In Section \ref{sec:main} we revisit implicit Runge-Kutta methods and discuss the main idea behind solving the nonlinear system using an explicit stabilised algorithm. In addition, we discuss a specific application of the new explicit implementation of implicit methods in the context of advection-diffusion-reaction equations and present an algorithm for doing so. {A detailed convergence analysis {in the case of pure linear diffusion and linear diffusion-advection} is presented in Section \ref{sec:conv}}. %{\color{red}(exSDIRK4: I am not sure what the acronym stands for but it would be good to define it somewhere in the introduction}.
Finally, in Section \ref{sec:num} we present a number of numerical computations that illustrate the theoretical convergence analysis in the pure linear diffusion case, and the linear diffusion-advection case. {In addition, applying the new explicit stabilised implementation exSDIRK4 of an order 4 SDIRK method to several nonlinear advection-diffusion-reaction problems illustrates the wide applicability and versatility of our approach. }

%\begin{itemize}
%\item mention ADIM methods  (again the attempt is to reduce the cost of solving an optimization problem)
%\end{itemize}


%Selling points ways to structure the introduction
%\begin{itemize}
%\item  What is the bottleneck in implementing implicit methods? We offer an implementation that does not suffer from this bottleneck
%\item Why is this important? a) We can have an actual cost estimate for specific problems? In theory not competitive with "actual second order methods" but these methods are only theoretical for high dimensional PDEs. We do everything using linear algebra and almost no large matrix inversion b) This unlocks the potential of using high-order methods for solving PDEs 
%\end{itemize}




        %\section{L-stability and stabilised Runge-Kutta methods}
        \section{Preliminaries}
\label{sec:pre}
In this paper, we are interested in the numerical integration of systems of ODEs of the form 
\begin{equation}\label{eq:ode}
		\dot y(t)=g(y(t)),\qquad y(0)=y_0, \quad 
t>0	\end{equation}
where $y(t) \in \mathbb{R}^{d}$ and the vector field $g: \mathbb{R}^{d} \mapsto \mathbb{R}^{d}$ is assumed smooth.  In this section, we recall the main tools for stability and accuracy analysis of Runge-Kutta methods.

%{\color{red}\begin{itemize}
%\item Some generic facts about Runge-Kutta methods (Butcher tableu, difference between implcit and explicit and in there mention potential benefits of diagonally implicit methods)
%\item Partitioned Runge-Kutta methods (discussed them in the context of equation (16) advection-diffusion-reaction equation) in there maybe discuss briefly about order conditons etc
%\end{itemize}}
        
        
        \subsection{Stability function and stability domain}
        
        Let $\lambda\in \IC $ and consider the Dahlquist test equation
        \begin{equation}\label{eq:testode}
            \dot y(t) = \lambda y(t), \quad y(0)=y_0\in \R.
        \end{equation}
        In the case where $\lambda \in \IC^-=\{z \in \IC;\,\text{Re}(z) \leq 0\}$, equation \eqref{eq:testode} has a bounded solution for  $t \rightarrow +\infty$. The simplest Runge-Kutta method for solving \eqref{eq:ode} is the explicit Euler method which when applied to \eqref{eq:testode} yields the recurrence
        \begin{equation}
            y_{n+1}=(1+\lambda \Delta t)y_{n}
        \end{equation}
        More generally, a Runge-Kutta method with step size $\Delta t$ applied to \eqref{eq:testode} produces the following recurrence
        \begin{equation} \label{eq:rec_dal}
        y_{n}=R(\Delta t\lambda)y_{n-1}=R(\Delta t \lambda)^ny_0,
        \end{equation}
        where $R(z)$ is called the stability function and the stability domain is defined by
        \begin{equation*}
            \calS=\{z\in\IC\,;\,|R(z)|\leq 1\}.
        \end{equation*}
        Note that $z \in \calS$ ensures that $y_{n}$ in \eqref{eq:rec_dal} remains bounded as $n \rightarrow \infty$ similarly to the true solution. In the case of the explicit Euler we have that 
        $
        \mathcal{S}_{\text{E}}=\{|z+1|\leq 1 \}
        $
        which corresponds to a disc of radius one, centred at $-1$ on the complex plane. For stiff problems where $|\lambda| \gg 1$ this implies a severe time-step restriction namely $\Delta t < 2/|\lambda|$ when $\lambda <0$. In particular, for spatial discretization of the heat equation eigenmodes are typically of size $\mathcal{O}(\Delta x ^{-2})$ which would in turn imply the CFL condition $\Delta t  < c \Delta x^{2}$.

        To avoid such severe time-step restrictions one solution is to use implicit methods. The simplest one of those is the implicit Euler method which when applied to  \eqref{eq:testode} yields the recurrence
        \[
        y_{n+1}=\frac{1}{1-\lambda \Delta t} y_{n}.
        \]
        In this  case we have $\mathcal{S}_{\text{IE}}=\{|z-1| \geq 1 \}$ which contains $\mathbb{C}^{-}$. This property is known as A-stability and can only hold in the case of implicit methods.  A-stability means that there is no restriction on the time-step size for eigenvalues with negative real parts contrary to the {explicit} Euler method and other explicit methods. Nevertheless, in the case of severely stiff  dissipative problems, a desirable additional property to A-stability is $R(\infty)=0$ which is known as L-stability \cite{HaW96}.      
        
        
        
       % it might still be the case that the stability function might  very close to $1$ when approaching infinity which leads to very slow convergence \cite[Chap.\ts4]{HaW96} 
        
        %problems, since for such methods the stability function might be very close to $1$ when approaching infinity which leads to very slow convergence \cite[Chap.\ts4]{HaW96}. In such cases, we prefer Runge-Kutta methods with the following property

       \subsection{Higher order Implicit Runge-Kutta methods}
       \label{subsec:hoirk}
        The implicit Euler is an example of an $L$-stable integrator meaning that is particularly suitable for integrating stiff equations. However, from the point of view of accuracy, it is a first-order method. To address this issue families of $L$-stable {Runge-Kutta methods have been proposed \cite{HaW96}} with an arbitrarily higher order of accuracy. Important examples are the families of the Radau and Lobatto Runge-Kutta methods are that achieve orders $2m-1$ and $2m-2$ respectively for any number $m$ of {internal stages \cite{I96}}. In particular, Radau with $m=1$ reduces to the implicit Euler method, while in the case where $m=2$  we obtain {the following implicit Runge-Kutta method for integrating \eqref{eq:ode},}
        \begin{subequations} \label{eq:radau2}
        \begin{eqnarray} 
         Y_{1}&=& y_{n}+\frac{5\Delta t}{12 }g\left(Y_{1}\right)-\frac{\Delta t}{12}g(Y_{2}), \\
       Y_{2}&=&  y_{n}+\frac{3\Delta t}{4 }g\left(Y_{1}\right)+\frac{\Delta t}{4}g(Y_{2}), \\
       y_{n+1}&=& Y_{2}.
       \end{eqnarray}
       \end{subequations}
  For a general $m$-stage implicit Runge-Kutta method with coefficients $a_{ij},b_j, i,j=1,\ldots,m$, given the numerical solution $y_n$ at time $t_{n}$  one needs to solve the following nonlinear system in dimension $m\cdot d$ with unknowns the internal stages $Y_i\in\R^d,i=1,\ldots,m$,
        \begin{equation}\label{eq:nonlinRK}
            Y_i= y_n + \Delta t \sum_{j=1}^m a_{ij} g(Y_j),\quad i=1,\ldots,m,
        \end{equation}
        to obtain the solution  $y_{n+1}$ at time $t_{n+1}$ by 
        \begin{equation} \label{eq:nonlinRKfinal}
            y_{n+1}= y_n + \Delta t \sum_{i=1}^m b_i g(Y_i).
        \end{equation}
        It is common to represent Runge-Kutta methods using a Butcher tableau
        \begin{equation}
            \begin{array}{c|c}
                 c&\calA\\
                 \hline
                 & b
            \end{array}
        \end{equation}
        where $\calA$ is the matrix of coefficients $\{a_{ij}\}_{i,j=1}^{m}$, $b$ is the weights vector $\{b_i\}_{i=1}^{m}$, and the vector $c$ is such that $c_i=\sum_{j=1}^{m} a_{ij}$.
        In compact matrix and tensor notations, \eqref{eq:nonlinRK} can be written as
        %\textcolor{red}{\big($I_d$ instead of $\mathbbm{1}_d$, also in other locations!!!\big)}
\begin{equation}\label{eq:nonlinearimp}
	\boldsymbol{Y}
	=
	\mathbbm{1}_m \otimes y_n
	+
	\Delta t (\calA \otimes I_d)
	g[\boldsymbol{Y}]
	\end{equation}
where $\calA \in \R^{m\times m}$ is the Runge-Kutta matrix of coefficients $a_{ij}$ and 
\begin{equation} \label{eq:summary}
\mathbbm{1}_m = \begin{pmatrix}
	1\\
	1\\
	\vdots\\
	1\\
	\end{pmatrix}\in\R^m,\quad
\boldsymbol{Y}=\begin{pmatrix}
	Y_1\\
	Y_2\\
	\vdots\\
	Y_m\\
	\end{pmatrix}\in\R^{m \cdot d},\quad
 g[\boldsymbol{Y}] = \begin{pmatrix}
	g(Y_1)\\
	g(Y_2)\\
	\vdots\\
	g(Y_m)\\
	\end{pmatrix}\in\R^{m \cdot d}
 \end{equation}
         In order to implement a general $m$-stage  implicit  Runge-Kutta method, for example, an $m$-stage Radau method, one employs Newton's method or variants of it. These methods make use of the Jacobian of the system which is of size $m d\times md$. Diagonally implicit Runge-Kutta methods (DIRK) are an attractive alternative to reduce this computational cost. This is because in this case, the matrix $\calA$ in \eqref{eq:nonlinearimp} is lower triangular, which implies that the Jacobians used by the Newton method or its variants to solve \eqref{eq:nonlinearimp} are of dimension $d\times d$. Furthermore, when all the diagonal coefficients are equal, the method is referred to as singly diagonally implicit (SDIRK), which may offer the potential to further reduce computational costs \cite{HaW96}. A particular diagonal implicit method of interest in this paper is SDIRK4 as presented in \cite[Table IV.6.5, p. 100]{HaW96} and whose Butcher tableau is given by
         \begin{equation}\label{eq:sdirk4}
\begin{array}
{c|ccccc}
\\
 & \frac14\\[1mm]
 &\frac12 &\frac14 \\[1mm]
& \frac{17}{50}& \frac{-1}{25}&\frac14 \\[1mm]
& \frac{371}{1360} &\frac{-137}{2720} &\frac{15}{544} &\frac14\\[1mm]
&\frac{25}{24} &\frac{-49}{48} &\frac{125}{16} &\frac{-85}{12} &\frac14\\[1mm]
\hline
y_1&\frac{25}{24} &\frac{-49}{48} &\frac{125}{16} &\frac{-85}{12} &\frac14\\[1mm]
\hat y_1&\frac{59}{48} &\frac{-17}{96} &\frac{225}{32} &\frac{-85}{12} &0\\[1mm]
err&-\frac{3}{16} &\frac{-27}{32} &\frac{25}{32} & 0 &\frac14
\end{array}.
\end{equation}
\begin{remark}%
Note that for SDIRK4 we have $y_{n+1}=Y_m$, where $m=5$. We call such methods stiffly accurate \cite{HaW96}. 
{The coefficient $\gamma=1/4$ is chosen in \cite[Sect.\thinspace IV.6, p. 99]{HaW96} to achieve both $L$-stability and a favorable error constant with an embedded method for error control and a continuous output of order 3. In \cite[Sect.\thinspace IV.6, p. 99]{HaW96}, a numerical better choice with $\gamma=4/15$ and numerical values of the $L$-stable Runge-Kutta coefficients is also proposed but we choose in the paper the SDIRK4 method \eqref{eq:sdirk4} for simplicity due to its nice rational coefficients.}
\end{remark}
\begin{remark} \label{rem:variable_step}
{
The SDIRK4 method in \eqref{eq:sdirk4} also includes an embedded method $\hat y_1$ for a variable time-step strategy \cite{HaW96} based on the error estimator 
\begin{equation} \label{eq:shampine}
err=J^{-1}(y_1-\hat y_1)
\end{equation}
where each new step $y_1$ with time step $\Delta t$ is accepted if $err$ is small enough compared to the tolerance $tol$ or rejected otherwise. The next time step is then computed as $\Delta t_{new} = \xi \Delta t (tol/err)^{1/5}$ where $\xi=0.8$ is a safety factor and $tol$ is the tolerance prescribed by the user.
The term $J^{-1}$ in \eqref{eq:shampine} denotes the inverse of the Jacobian matrix $J=I_d-\Delta t \gamma \frac{\partial g}{\partial y}$ computed during the Newton method or its variant at iterates of the internal stages $Y_i$, and hence $J^{-1}$ is already computationally available with minor overhead in practice. This is a standard methodology for stiff problems proposed by Shampine, as presented in \cite[Sect.\thinspace IV.8]{HaW96}, to circumvent the lack of $L$-stability of the embedded method and hence  stabilize the error estimator which would otherwise be too pessimistic and fail for stiff problems. 
}
\end{remark}
         %\begin{subequations}
         %\begin{eqnarray}
          %Y_{1} &=& \\
          %Y_{2} &=& \\
          %Y_{3} &=&  \\
          %Y_{4} &=&  
         %\end{eqnarray} 
         %\end{subequations}
        %with 
        %\[
        %y_{n+1}=y_{n}
        %\]
       
        \subsection{Explicit stabilised Runge-Kutta methods} \label{sec:rkc}
        Even with the dimensionality reduction of diagonally implicit Runge-Kutta methods in terms of the Jacobians used for solving the nonlinear equations, implicit methods in high dimensions become prohibitively expensive. On the other hand, explicit schemes avoid the solution of such high-dimensional systems, but have bounded stability domains.  For this reason, standard explicit methods face severe step size restrictions when integrating stiff problems. Explicit stabilised Runge-Kutta methods represent a good compromise, they are fully explicit schemes with extended stability domains over the negative real axis. 


        
        %Explicit stabilised methods were originally constructed to solve diffusion problems with eigenvalues in a long strip around the negative real axis \cite{SSV98,AbM01,Abd02}. Later, more elaborated schemes of this type were introduced to solve a wide variety of stiff problems ranging from advection-diffusion-reaction PDEs to SDEs and optimal control problems \cite{AbL08,Zb11,AVZ13,AV13,AAV18,AV21,Alm21}. 
        
        Explicit stabilised Runge-Kutta methods are constructed to have (nearly) maximal stability domain sizes along the negative real axis constrained by the fact that they remain explicit. In particular,  for an $s$-stage method the stability polynomial would be of degree $s$. The goal is then  to find, for a given integer $p$ (the order of the method), a polynomial of degree $s$ (the number of internal stages of the method) satisfying 
        \begin{equation*}
            R_s(z)=1+z+\frac{z^2}2+\dots+\frac{z^p}{p!}+\alpha_{p+1}z^{p+1}+\dots+\alpha_sz^s,
        \end{equation*}
        such that $|R_s(z)|\leq1$ for $z\in [-L_s,0]$ with $L_s>0$ is as large as possible. For $p=1$, the optimal value of $L_s$ is $2s^2$ \cite{Abdrev,HaW96}. The corresponding polynomial is the shifted Chebyshev polynomial $R_s(z)=T_s(1+z/s^2)$ where $T_s(\cos\theta) = \cos(s\theta)$ is the first kind Chebyshev polynomial of degree s. This polynomial oscillates between $-1$ and $1$ which introduces singularities in the stability domain and prevents convergence for some eigenvalues. For this reason, we modify the stability function in the following way
        \begin{equation}\label{eq:stabcheb}
            R_s(z)=\frac{T_s(\omega_0+\omega_1z)}{T_s(\omega_0)},
        \end{equation}
        where $\omega_0=1+\eta/s^2$, $\omega_1=T_s(\omega_0)/T_s'(\omega_0)$, and $\eta$ is called the damping parameter. The polynomial \eqref{eq:stabcheb} now oscillates between $-\alpha_s(\eta)$ and $\alpha_s(\eta)$ for every $z\in[-L_{s,\eta},-\delta_\eta]$, where
        \begin{equation}\label{eq:alphaL}
            \alpha_s(\eta)=\frac1{T_s(\omega_0)}<1, \quad L_{s,\eta}=\frac{1+\omega_0}{\omega_1},
        \end{equation}
        where $\delta_\eta$ is a positive real close to zero ($R_s(0)=1$). %For diffusion dominated problems, a small damping $\eta=0.05$ is usually sufficient. The stability region grows in the imaginary direction as the damping parameter increases, which turns out to be very useful in the context of SDEs and PDEs with large advection terms \cite{AbL08,AV13,Alm21}. Using large values of $\eta$ was also exploited in the context of optimization in \cite{EVVZ21} to speed up convergence. We will discuss this idea in details in the next section.

        Using the well-known three-term recurrence relation of the first kind Chebyshev polynomials $T_{j+1}(z)=2zT_j(z)-T_{j-1}(z)$, for $j\geq2$ (with $T_0(z)=1$ and $T_1(z)=z$), we can derive an efficient and stable implementation (in terms of memory and rounding errors) of the so-called Chebyshev method or the first order Runge-Kutta-Chebyshev (RKC) method with stability function \eqref{eq:stabcheb},  as analyzed by \cite{vanderHouwen1980} (also for second order explicit stabilised methods). When applied to non-linear problems where the Jacobian of $g$ has negative real eigenvalues, this method allows for taking much larger time steps than standard explicit methods would allow. Algorithm \ref{alg:rkc} summarises one step of the first order $s$-stage RKC method with step-size $\Delta t$ for the solution of \eqref{eq:ode}. 
\begin{algorithm}  
\caption{{First order} Runge-Kutta Chebyshev for the time integration of \eqref{eq:ode}} 
\textbf{Input:} Initial state $y_{n}$, damping $\eta$, upper bound $L$ for the eigenvalues of  the Jacobian of $g$ at $y_{n}$, time step $\Delta t$
%\vspace{5pt}        
        \begin{equation}\label{eq:rkc1}
            \begin{split}
            K_0 &= y_n,\quad K_1=K_0+\mu_1 \Delta t g(K_0),\\
            K_j &= \mu_j \Delta t g(K_{j-1})+\nu_jK_{j-1}+(1-\nu_j)K_{j-2},\quad j=2,\dots,s,\\
            y_{n+1} &= K_s,
            \end{split}
        \end{equation}
        where
        $$\mu_j=\frac{2\omega_1T_{j-1}(\omega_0)}{T_j(\omega_0)},\quad \nu_j=\frac{2\omega_0T_{j-1}(\omega_0)}{T_j(\omega_0)}$$
        and $s$ should be chosen large enough such that $L_{s,\eta}>L\Delta t$ 
        \newline
        \textbf{Output:} Solution $y_{n+1}$ after a time step of size $\Delta t$ 
        \label{alg:rkc}
\end{algorithm}
Note that there are three parameters needed for implementing this Algorithm, the damping $\eta$, the number of stages $s$ as well as the  step-size $\Delta t$. 
{For a fixed value of the damping parameter $\eta$ and a given time step $\Delta t$ one needs to choose the stage parameter $s$ in such a way that $L_{s,\eta} \geq L \Delta t$.
Typically for ODE integration one takes $\eta$ small (e.g $\eta=0.05$) which gives a nearly optimal stability domain size $L_{s,\eta} \geq (2-\frac{4}{3}\eta)s^{2}$, and a stage parameter
\[
s=\left\lceil \sqrt{\frac{L \Delta t}{2-\frac{4}{3}\eta}}\right\rceil
\] where $L$ is an upper bound for the largest eigenvalue of the Jacobian of $g$ evaluated at $y_{n}$.
For larger values of damping parameter $\eta$, the width $L_{s,\eta}$ of the stability domain reduces further but still grows quadratically with $s$ and the number of stages $s$ is then chosen large enough accordingly, see  for instance \cite[Remark 6.1]{AAV18} for details.
%$L_{s,\eta} \geq 2\Omega(\eta)^{-1} s^2$ where $\Omega(\eta)^{-1} =\tanh(\sqrt{2\eta})/\sqrt{2\eta}$, see for instance \cite[Remark 6.1]{AAV18}, and the number of stages $s$ is then chosen as 
%$
%s=\left\lceil \sqrt{L \Delta t\Omega(\eta)^{-1}/2}\right\rceil.
%$
}
\subsection{Application to optimization}


In the recent paper \cite{EVVZ21} the authors utilised \eqref{eq:rkc1} to solve efficiently stiff optimization problems by an algorithm called  Runge-Kutta Chebyshev descent {(RKCD)}.  Here, by stiff optimization problems, we mean the case where the Hessian of the objective function has large eigenvalues distributed in the positive half-plane close to the real axis. In particular, consider the problem 
	\begin{equation}\label{eq:min}
		\min_{x\in\ \R^d}f(x),
	\end{equation}
	where the dimension $d$ is large and $f:\R^d\to\R$,  is $\ell$-strongly convex differentiable function, with $L$-Lipschitz derivative, and such that the real parts of the eigenvalues of its Hessian matrix $\mathcal{H}f(x)$ are strictly positive for all $x\in\R^d$. The condition number of the Hessian is bounded by $\kappa=\frac{L}{\ell}$.
 
 A starting point for solving the minimization problem \eqref{eq:min} is to consider discretizations of the gradient flow
	\begin{equation}\label{eq:gf}
		\dot x(t) = -\nabla f(x(t)), \qquad x(0)=x_0 \in \R^d,
	\end{equation}
	where $x_0$ is an initialization. Applying an explicit Euler method with step size\footnote{We will use $\Delta t$ for the actual time step of numerical integrators for differential equations and $h$ for optimization steps.} $h$ to \eqref{eq:gf} is equivalent to the gradient descent method
	\begin{equation}\label{eq:gradient descent}
		x_{k+1}=x_k-h\nabla f(x_k),\qquad k=0,1,2,...
	\end{equation}
	where $x_k$\footnote{When we describe an optimization algorithm to approximate a minimizer $x_{*}$ we will use the notation $x_{k}$ for the $k$-th iterate of the algorithm.} is the numerical approximation of $x(t_k)$. In the strongly convex setting described above when  ($L\gg l$), this method faces a severe restriction on $h$ to be stable. Furthermore, it is known that the optimal $h$ in terms of convergence speed towards the minimizer is given by {$h=\frac{2}{L+\ell}$ \cite{N18}}.  This implies that the number of iterations needed to find the minimizer of $f$ scales linearly with the condition number $\kappa$. Note that this behaviour is suboptimal since the most efficient optimization methods that make use only of the gradient of $f$ need $\mathcal{O}(\sqrt{\kappa})$ iterations to {find the minimum \cite{N18}}.  

To improve the convergence speed of the explicit Euler method one can apply a more sophisticated discretization of the gradient flow \eqref{eq:gf}. For example, one can use the first order RKC \eqref{eq:rkc1} to integrate  \eqref{eq:gf}. In particular, if one has estimates for the smallest and largest eigenvalues of the Hessian matrix $\nabla^{2} f$, then by choosing the parameters $\eta, h,$ and $s$ in \eqref{eq:rkc1} appropriately one can obtain a very efficient optimization method called RKCD \cite{EVVZ21} for strongly convex problems, see also Algorithm \ref{alg:rkcd}. In the case of RKCD, a calculation of $x_{k+1}$ for $x_{k}$ requires $s$ gradient evaluation. Now the following proposition from \cite{EVVZ21} characterizes the speed of convergence of iterates $x_{k}$ towards the minimizer  $x_{\star}$ when $f$ is quadratic in \eqref{eq:min}.

\begin{algorithm} [t]
\caption{{\cite{EVVZ21}} Runge-Kutta Chebyshev descent method (RKCD)} 
%\vspace{5pt}


\textbf{Input:}
%\begin{itemize}
%\item 
The gradient force $g=-\nabla f$ of a differentiable function $f:\mathbb{R}^d\rightarrow\mathbb{R}$.
%\item 
Damping $\eta$. Lower and upper bounds $\ell$ and $L$ for eigenvalues of the Jacobian of $g$.  Initialisation $x_0\in\mathbb{R}^d$. 
%\end{itemize}
%\vspace{5pt}




\textbf{Body:} 
%\begin{itemize}
%\item 
%Until convergence, for $k\in \{0,1,2,3,\ldots\}$, \textbf{repeat}:
\begin{itemize}
  \item Compute the time step $h$ (learning rate) and the number of stages $s$ using 
  \begin{equation}  \label{eq:defn of parameters1}
  s=\left\lceil\sqrt{(\kappa-1)\eta/2}\right\rceil,\quad h=\frac{\omega_0-1}{\omega_1\ell}
  \end{equation}
  where $\kappa=L/\ell$  and
\begin{equation}\label{eq:defn of parameters}
\omega_{0}=1+\frac{\eta}{s^{2}}, \quad \omega_{1}
=\frac{T_{s}(\omega_{0})}{T'_{s}(\omega_{0})},
\end{equation}
with $T_s$ is the Chebyshev polynomial of the first kind with degree $s$. 
%\item {Compute $h$ and $s$ using~\eqref{eq:condition_2a},\eqref{eq:defn of cnd no} where $L,\mu$ are upper and lower bounds for the eigenvalues of the Hessian $\nabla^2 f(x_n)$.}
\item Until convergence, for $k\in \{0,1,2,3,\ldots\}$, \textbf{repeat}:
\begin{itemize}
\item Set $x_k^0 = x_k$ and $x_k^1 = x_k^0 + h\mu_1 g(x_k)$ with $\mu_1 = \omega_1/ \omega_0$
\item For $j\in \{2,\cdots,s\}$, \textbf{repeat}:
%\begin{itemize}
%\item 
%For every $j\in \{2,\cdots,s\}$, set 
$$x_k^j = \mu_j h g(x_k^{j-1}) + \nu_j x_{k}^{j-1}-(\nu_j-1) x_{k}^{j-2},$$ 
% where 
\[
\mu_j = \frac{2\omega_1 T_{j-1}(\omega_0)}{T_j(\omega_0)},
\quad \nu_j = \frac{2\omega_0 T_{j-1}(\omega_0)}{T_j(\omega_0)}.
\]
%\end{itemize}
\item Set $x_{k+1}=x_{k}^s$. 
\end{itemize}
\end{itemize}
%\vspace{5pt}

\vspace{0.1cm}

\textbf{Output: } 
%\begin{itemize}
%\item 
Estimate $\widehat{x}={x}_{n+1}$ of a minimiser of  ~\eqref{eq:min}. 
%\end{itemize}

\label{alg:rkcd}
\end{algorithm}



 
 

 
 \begin{proposition} \label{prop:quad_RKC}
 {\cite{EVVZ21}.} {Let $f(x)=\frac12 x^TCx+b^Tx+c$ and consider the minimization problem \eqref{eq:min}}. Now let $x_{k}$ be the iterate of Algorithm \ref{alg:rkcd} with parameters chosen according to \eqref{eq:defn of parameters1} and
 \eqref{eq:defn of parameters}. 
Then the following {convergence estimates hold,}
 \begin{equation} \label{eq:ratte}
         \|x_{k}-x_*\|\leq \alpha^{k}_{s}(\eta)\|x_0-x_*\|,
     \end{equation}
     and
 \begin{equation}
         f(x_{k+1})-f(x_*)\leq \alpha^{2}_{s}(\eta)(f(x_k)-f(x_*)),
     \end{equation}
where  $x_{*}$ is the unique minimizer of $f$, and $\alpha_{s}(\eta)$ satisfies \eqref{eq:alphaL}. 

\end{proposition}

%\begin{remark}
    A straightforward calculation shows that the learning rate $h$ in \eqref{eq:defn of parameters1} satisfies
    $
    h= \tfrac{\eta}{\ell} + \bigo(\eta^2 + s^{-2}),
    $
    for $\eta\rightarrow 0, s\rightarrow +\infty$ which makes it mildly dependent on $s$. %Furthermore, it is possible to show that 
    %\begin{equation} \label{eq:h_bound}
    %h\leq \frac{\eta}{\ell}.
    %\end{equation}
%\end{remark} 
Precisely, the following estimates holds.
\begin{lemma}
The learning rate $h$ defined in \eqref{eq:defn of parameters1} satisfies the bound
\begin{equation} \label{eq:h_bound}
h\leq \frac{\eta e^\eta}{\ell}.
   \end{equation}
\end{lemma}
\begin{proof}
We recall the following formula for all $s\geq 1,k\geq 0$,
\begin{equation} \label{eq:Tsprime}
T_s^{(k)}(1) = \prod_{j=0}^{k-1} \frac{s^2-j^2}{2j+1} \leq s^{2k}.
\end{equation} 
Using $1/\omega_1 = T_s'(\omega_0)/T_s(\omega_0)\leq T_s'(\omega_0)$, and performing a Taylor expansion of $T_s'(\omega_0)$ for $\eta$ close to 
$0$ yields
$$
\ell h=\frac{\eta}{s^2\omega_1} \leq \frac{\eta T_s'(\omega_0)}{s^2}
=\sum_{k=0}^{s-1} \frac{T_s^{(k+1)}(1)}{k! s^{2k+2} }{\eta^{k+1}} \leq 
\sum_{k=0}^{s-1} \frac{\eta^{k+1}}{k!} \leq \eta e^\eta
$$
where we used \eqref{eq:Tsprime} in the last inequality, which concludes the proof. 
\end{proof}
Using this proposition and the properties of $\alpha_{s}(\eta)$  it is possible to show that the effective rate of convergence of function iterates towards the minimizer $x_{*}$ scales like $1-c(\eta)/\sqrt{\kappa}$  when $\kappa \gg 1$ \cite{EVVZ21}, exhibiting thus accelerated behavior from an optimization point of view. This, in turn, implies that  to calculate $x_{\star}$ with target  accuracy $\epsilon$ one needs to do $\mathcal{O}(\sqrt{\kappa}\log{\epsilon^{-1}})$ gradient evaluations. 

%Algorithm \ref{alg:rkcd} can also be used for ....

%In addition, a modification of \eqref{eq:rkc1} was proposed in \cite{EVVZ21} to deal with the minimization problem \eqref{eq:min} for functions $f$ that are quadratic plus a small nonquadratic perturbation. 
 
%\subsection{Methods for solving ARD PDEs} 
        
	\section{A new explicit implementation of SDIRK  methods}
 \label{sec:main}
%	We aim to implement excellent well-established high-order implicit Runge-Kutta methods explicitly by combining some optimization techniques with an explicit stabilised method of order $1$. We will explain the methodology in the rest of the current section.
%We now discuss 
%\[
%\dot{y}= F_{D}(y)+F_{A}(y)+F_{R}(y)
%\]    
%write down the system with $G_{D}, G_{A}, G_{R}$	
 
%To deal with stiff gradient flows \eqref{eq:gf} and avoid time-step size restriction, another approach is to discretize \eqref{eq:gf} using an implicit Euler method with a step size $h$
%	\begin{equation}\label{eq:ie}
%		x_{n+1}=x_n-h\nabla f(x_{n+1}),\qquad n=0,1,2,...
%	\end{equation}
%	It is easy to verify that $x_{n+1}$ in \eqref{eq:ie} is the unique minimizer of the proximal problem
%	\begin{equation}\label{eq:prox}
%		\min_{x\in\ \R^d}hf(x)+\frac12 \|x-x_n\|_2^2.
%	\end{equation}
%	It is reasonable to think of applying a gradient descent method with step size $h$ to the gradient flow of \eqref{eq:prox}
 %\begin{equation}\label{eq:gfprox}
 %    \dot x=-h\nabla f -x + x_n
% \end{equation}
% to avoid implicitness and reduce the step size restriction when $h<1$, but this will not reduce much the restriction on $h$ especially if $h$ is large. This pushes us to think of applying an explicit stabilised method to \eqref{eq:gfprox} with a step size $h=\bigo(1)$ and a suitable number of stages $s$, and iterate until convergence. The big advantage of this approach is that we will obtain the solution of an L-stable implicit method without solving any linear and/or nonlinear systems. We aim to apply this idea to powerful high-order implicit methods with perfect stability properties.
As discussed in Section \ref{subsec:hoirk}, implicit methods are usually implemented using Newton's method or its variants which can be memory demanding and delicate for ill-conditioned problems in particular for high-dimensional problems arising from discretizing PDEs.  Here, we propose an alternative way of implementing implicit Runge-Kutta methods 
\eqref{eq:nonlinRK}-\eqref{eq:nonlinRKfinal}
based on explicit stabilised methods. To do this, one needs to find the solution $\boldsymbol{Y}_{n+1}$ to the system of non-linear equations \eqref{eq:nonlinearimp}  given the current state $y_{n}$. Our starting point for achieving this is 
to introduce the auxiliary system of ODEs,
\begin{equation}\label{eq:gf1}
			\dot {\boldsymbol{Y}}(t)=-\boldsymbol{Y}(t)
	+\mathbbm{1}_m \otimes y_n
	+
	\Delta t (\calA \otimes I_d)
	g[\boldsymbol{Y}(t)]
\end{equation}  
that satisfies under appropriate assumptions 
\[
\lim_{t  \rightarrow +\infty} \boldsymbol{Y}(t) = \boldsymbol{Y}_{n+1}.
\]
We have thus reformulated the solution of an $m$-stage implicit Runge-Kutta method at time $t_{n+1}$ as the steady-state solution of the  ODE \eqref{eq:gf1}. 
In the simplest case of the implicit Euler method and assuming a gradient structure $g(y)=-\nabla V(y)$ in \eqref{eq:ode}, it is easy to see that  \eqref{eq:gf1} also has a gradient structure \eqref{eq:gf} with 
%\begin{equation} \label{eq:f_new}
$ f(y)=\Delta t V(y)+\frac12 \| y - y_n\|_2^2. $ 
%\end{equation}
In this simple situation, following the ideas in \cite{EVVZ21} one can apply RKCD directly to compute $\boldsymbol{Y}_{n+1}$. 
However, this gradient structure assumption is not satisfied in general for a diagonally implicit Runge-Kutta method, even when \eqref{eq:main} has a gradient structure. In particular, \eqref{eq:gf1} in this case becomes
\begin{equation} \label{eq:gff}
\dot {\boldsymbol{Y}}(t)=-\boldsymbol{Y}(t) +\mathbbm{1}_m \otimes y_n
	+
	\Delta t (\calA \otimes I_d)(F_{D}[{\boldsymbol{Y}}(t)]+F_{A}[{\boldsymbol{Y}}(t)]+F_{R}[{\boldsymbol{Y}}(t)])
\end{equation}
As we will see, without this gradient assumption, applying Algorithm \ref{alg:rkcd}  to \eqref{eq:gff} would lead to arbitrarily slow convergence to $\boldsymbol{Y}_{n+1}$ as the stiffness grows. To tackle this issue, and inspired by the partitioning $\calA=\gamma I_m+(\calA-\gamma I_m)$ of the Runge-Kutta coefficients into diagonal part and strictly lower diagonal part proposed in \cite {VSC92},
we consider the following partition of \eqref{eq:gff},
\begin{subequations} \label{eq:gff_part}
\begin{align}
\dot{{\boldsymbol{Y}}}(t) \;=\;& \mathbbm{1}_m \otimes y_n - {\boldsymbol{Y}}(t) \;+\; \Delta t\, \gamma\, I_m \otimes I_d \, F_D[ {\boldsymbol{Y}}(t)] &&=:~G_D( {\boldsymbol{Y}}(t))\\
&\quad +~ \Delta t\, (\calA-\gamma I_m)\otimes I_d \, F_D[ {\boldsymbol{Y}}(t)] \;+\; \Delta t\, \calA\otimes I_d \, F_A[{\boldsymbol{Y}}(t)] &&=:~G_A( {\boldsymbol{Y}}(t))\\
&\quad +~ \Delta t\,\calA\otimes I_d \, F_R[ {\boldsymbol{Y}}(t)] &&=:~G_R( {\boldsymbol{Y}}(t)),
\end{align}
\end{subequations}
where we split the coefficient matrix $\mathcal{A}$ of the
singly diagonally implicit Runge-Kutta method
into its diagonal part $\gamma I_m$ and its strictly lower triangular part $\mathcal{A}-\gamma I_m$. 
The new Algorithm \ref{alg:3} is then defined as follows.
\begin{algorithm} 
\caption{{New} Runge-Kutta Chebyshev implementation of diagonally implicit Runge-Kutta methods \label{alg:3} }
%\vspace{5pt}


\textbf{Input:}
%\begin{itemize}
%\item 
The vector fields $\mathbb{R}^d\rightarrow\mathbb{R}^d$, $F_D$ (diffusion terms), $F_A$ (advection terms), $F_R$ (reaction terms).
%\item 
The initial state $y_{n}\in\mathbb{R}^d$, an approximation ${\boldsymbol{x}}_0\in\mathbb{R}^{md}$ of ${\boldsymbol{Y}}_{n}$. Upper bounds $\lambda_{max}$ for the spectral radius of $\nabla F_D$.  
Diagonally Runge-Kutta coefficient matrix $A$ with value $\gamma>0$ on the diagonal, time-step size $\Delta t$. Damping $\eta$. 
%\end{itemize}
%\vspace{5pt}

\vspace{0.1cm}


\textbf{Body:} 
%\begin{itemize}
%\item 
%Until convergence, \textbf{repeat}:
\begin{itemize}
  \item Compute the time step $h$ (learning rate) and the number of stages $s$ according to 
  \[
  s=\left\lceil\sqrt{(\kappa-1)\eta/2}\right\rceil,\quad h=\frac{\omega_0-1}{\omega_1\ell}
  \]
  where $\ell=1$, $L=1+\gamma \Delta t \lambda_{max},\kappa=L/\ell$ and
\[
\omega_{0}=1+\frac{\eta}{s^{2}}, \quad \omega_{1}
=\frac{T_{s}(\omega_{0})}{T'_{s}(\omega_{0})},
\]
with $T_s$ is the Chebyshev polynomial of the first kind with degree $s$. 
%\item {Compute $h$ and $s$ using~\eqref{eq:condition_2a},\eqref{eq:defn of cnd no} where $L,\mu$ are upper and lower bounds for the eigenvalues of the Hessian $\nabla^2 f(x_n)$.}
\item Until convergence, for $k\in \{0,1,2,3,\ldots\}$, \textbf{repeat}:
\begin{itemize}
\item Define 
\begin{align*}
G({\boldsymbol{x}}):&=G_D({\boldsymbol{x}}) - G_D({\boldsymbol{x}}_k)
+ J^{-1}\!\bigl(G_D({\boldsymbol{x}}_k)+G_A({\boldsymbol{x}}_k)+G_R({\boldsymbol{x}}_k)\bigr),
\\
J &= I_{md} - h\, G'_R({\boldsymbol{x}}_k),
\end{align*}
%where
%\begin{align*}
%G_D({\boldsymbol{x}}) :=& \mathbbm{1}_m \otimes y_{n} - {\boldsymbol{x}} \;+\; \Delta t\, \gamma\, \mathbbm{1}_m \otimes \mathbbm{1}_d \, F_D({\boldsymbol{x}}) \\
%G_A({\boldsymbol{x}}) :=& \Delta t\, (A-\gamma %\mathbbm{1}_m)\otimes \mathbbm{1}_d \, F_D( {\boldsymbol{x}}) \;+\; \Delta t\, A\otimes \mathbbm{1}_d \, F_A( {\boldsymbol{x}}) \\
%G_R( {\boldsymbol{x}})  :=& \Delta t\, A\otimes \mathbbm{1}_d \, F_R( {\boldsymbol{x}}).
%\end{align*}
\item Set ${\boldsymbol{x}}_k^0 = {\boldsymbol{x}}_k$ and ${\boldsymbol{x}}_k^1 = {\boldsymbol{x}}_k^0+ h\mu_1 G({\boldsymbol{x}}_k)$ with $\mu_1 = \omega_1/ \omega_0$
\item For $j\in \{2,\cdots,s\}$, \textbf{repeat}:
%\begin{itemize}
%\item 
%For every $j\in \{2,\cdots,s\}$, set 
$${\boldsymbol{x}}_k^j = \mu_j h G({\boldsymbol{x}}_k^{j-1}) + \nu_j {\boldsymbol{x}}_{k}^{j-1}-(\nu_j-1) {\boldsymbol{x}}_{k}^{j-2},$$ 
% where 
\begin{equation} \label{eq:rkcd_iter}
\mu_j = \frac{2\omega_1 T_{j-1}(\omega_0)}{T_j(\omega_0)},
\quad \nu_j = \frac{2\omega_0 T_{j-1}(\omega_0)}{T_j(\omega_0)}.
\end{equation}
%\end{itemize}
\item Set ${\boldsymbol{x}}_{k+1}={\boldsymbol{x}}_{k}^s$. 
\end{itemize}
\end{itemize}
%\vspace{5pt}

\vspace{0.1cm}

\textbf{Output: } 
%\begin{itemize}
%\item 
Approximation ${\boldsymbol{x}}_{k+1}$ of the Runge-Kutta solution ${\boldsymbol{Y}}_{n+1}\in\mathbb{R}^{md}$ after one step. 
%\end{itemize}

\label{alg:explicitimplicit}
\end{algorithm}




\begin{remark}
Algorithm \ref{alg:explicitimplicit} 
calculates
 $\boldsymbol{Y}_{n+1}$ by applying  Algorithm \ref{alg:rkcd}  to the following new partitioned auxiliary problem:
\begin{align} \label{eq:part}
\dot{\boldsymbol{Y}}
&= G_D(\boldsymbol{Y}) - G_D(\boldsymbol{Y}_n)
+ J^{-1}\!\bigl(G_D(\boldsymbol{Y}_n)+G_A(\boldsymbol{Y}_n)+G_R(\boldsymbol{Y}_n)\bigr),
\\
J &= I_{md} - h\, G'_R(\boldsymbol{Y}_n).
\end{align}
Note that this partitioning preserves the steady state ${\boldsymbol{Y}}_{n+1}$ of \eqref{eq:gff} since 
\[
G_{D}({\boldsymbol{Y}}_{n+1})+G_{A}({\boldsymbol{Y}}_{n+1})+G_{R}({\boldsymbol{Y}}_{n+1})=0.
\]
\end{remark}

\begin{remark}
The partitioning \eqref{eq:gff_part} and the corresponding treatment of the $G_{R}$ term in \eqref{eq:part} assume that the reaction term $F_{R}$ is stiff. Note that the Jacobian inversion in this setting is cheap as it has a block diagonal structure thanks to the local nature of the reaction terms.  In practice, for a non-stiff reaction, one takes $G_{R}$ to be zero to avoid the Jacobian computations and then take \[G_{A}(\boldsymbol{Y}(t)):=\Delta t\, (\calA-\gamma I_m)\otimes I_d \, F_D[ {\boldsymbol{Y}}(t)] \;+\; \Delta t\, \calA\otimes I_d \, \left(F_A[{\boldsymbol{Y}}(t)]+ F_R[ {\boldsymbol{Y}}(t)]\right).
\]
This corresponds to treating together the advection and non-stiff reaction terms as first proposed for PRKC \cite{Zb11} and also used in PIROCK \cite{AV13}.
\end{remark}



In Algorithm \ref{alg:3}, we set $\ell=1$ for simplicity, as the spectral gap for the Jacobian of $G_D$. Note that if $\lambda_{\min}\geq 0$ denotes the smallest eigenvalue of the Jacobian of $-F_D$ one could use the more precise value $\ell=1+\Delta t\lambda_{\min}$. 
Concerning the choice of the damping parameter $\eta$, while the analysis applies for any value $\eta>0$, it was shown numerically in \cite{EVVZ21} that the value $\eta=1.17$ is an efficient choice for RKCD. For Algorithm \ref{alg:3}, we choose $\eta=4$ which reveals efficient in the considered numerical test problems (see Section \ref{sec:num}). Concerning the stopping criteria for the main iteration of Algorithm \ref{alg:3}, we use for simplicity the condition $\|\boldsymbol{x}_k-\boldsymbol{x}_{k-1}\|<tol_{it}$, with $tol_{it}=10^{-12}$. For the initialisation, we use for simplicity $\boldsymbol{x}_0=\mathbbm{1}_m\otimes y_n$ where $y_n$ is the current state.




%\begin{remark} \label{rem:gf}
%Consider the case where $g(y)=-\nabla V(y)$ in \eqref{eq:ode}. Then if the matrix $A$ is symmetric then the vector field  $(A \otimes \mathbbm{1}_d)
%		g(\boldsymbol{Y})$  can be written as $-\nabla  F( \boldsymbol{Y})$
%then \eqref{eq:gf1} corresponds to the gradient flow \eqref{eq:gf} for   
% \begin{equation} \label{eq:f_new}
%   f(\boldsymbol{Y})=\Delta t F(\boldsymbol{Y})+\frac12 \| \boldsymbol{Y} -\mathbbm{1}_m \otimes y_n\|_2^2.   
%  \end{equation}
%  Explain why this happens for implicit Euler
%\end{remark}



\section{Convergence and stability analysis} \label{sec:conv}
In this section, we prove the convergence of Algorithm \ref{alg:3} in two special cases, pure linear diffusion and linear advection-diffusion. Additionally, we provide the rationale behind the partition in equation \eqref{eq:gff_part} for a general advection-diffusion-reaction equation.

\subsection{The case of pure linear diffusion}
We  consider 
\begin{equation}\label{eq:heat}
\partial_t u=a\Delta u, \quad (t,x)\in[0,+\infty)\times(0,1)^{N}
\end{equation}
with periodic boundary conditions and a smooth initial condition $u(x,0)=1+e^{-50\norm{x-c}^2}$ where $c$ is the center of the spatial domain.  We  now  discretise in space using a finite difference method with a uniform mesh size $\Delta x$ in each spatial dimension.  We thus have 
\begin{equation} \label{eq:pure_diffusion}
F_{D}(y)= -Dy, \quad F_{A}(y)=0 \quad F_{R}(y)=0
\end{equation}
in \eqref{eq:main} with  $D \in \mathbb{R}^{d\times d}$ is the discrete Laplacian with periodic boundary conditions. In this case, we have that the bound for the spectral radius of $\nabla F_{D}(y)=-D$ is given by $\lambda_{\text{heat},max}=4 N a /\Delta x^{2}$, for $N=1,2,3$.

\paragraph{Implicit Euler}
We consider the solution of \eqref{eq:main} by implicit Euler when the different terms are given by \eqref{eq:pure_diffusion}. In this case, we have that $m=1$ and $\mathcal{A}=\gamma=1$ in \eqref{eq:gff} implying that 
\begin{equation} \label{pure_diffusion1}
G_{D}(y)= \mathbbm{1}_m\otimes y_{n}-y-\Delta t D y, \quad G_{A}(y)=0 \quad G_{R}(y)=0
\end{equation}
in \eqref{eq:gff_part}.  
We now use Algorithm \ref{alg:explicitimplicit} to calculate the solution of implicit Euler at time $t_{n+1}$ given the solution $y_{n}$ at time $t_{n}$. Note that in this case, because $G_{A}(y)=G_{R}(y)=0$  and $m=1$ \eqref{eq:gff} has a gradient structure and Algorithm \ref{alg:3} coincides with Algorithm \ref{alg:rkcd}. One would thus expect from Proposition \ref{prop:quad_RKC} that given $y_{n}$, the cost to approximate $y_{n+1}$ scales proportionally to the 
$\sqrt{\kappa}$ where $\kappa$ is the condition number of $G_{D}$ given by $\kappa=1+\gamma \Delta t \lambda_{\text{heat},max}=1+4Na/\Delta x^{2}$.  To illustrate this, we now vary $\kappa$ by changing the size of the mesh refinement $\Delta x= 1/4,1/8,1/16,1/32,1/64$ while keeping $\Delta t=5 \times 10^{-2}$. The results of this experiment are presented in Figure \ref{fig:3dLaplace} for $N=1$ and $N=3$ spatial dimensions, and show that the cost in terms of gradient evaluations indeed scales like $\mathcal{O}(\sqrt{\kappa})$ independently of the value of $N$. Here we have used a prescribed tolerance $\varepsilon=10^{-10}$ in the criteria for stopping the iterations.
%For the same values of $\Delta t$, we now fix $\Delta x=10^{-2}$ and plot in Figure \ref{fig:euler}b  the error $\|x_{k}-y_{1}\|^{2}$ as a function of the number of iterations $k$, where $y_1\simeq y(\Delta t)$ is calculated using Newton method up to machine precision. As we can see for both figures the results are as we expect according to Proposition \ref{prop:quad_RKC}. In both experiments, we choose $a=0.2$ and $u(x,0)=1+\sin{(2\pi x)}$. The damping parameter $\eta$ is set to $4$. 

%\begin{figure}
%    \centering
%    \begin{subfigure}{0.49\linewidth}
%    \centering
%    \includegraphics[width=\linewidth]{costvsdx.png}
%    \caption{Cost}
%    \end{subfigure}
%    \begin{subfigure}{0.49\linewidth}
%    \centering
%    \includegraphics[width=\linewidth]{increment_euler.png}
%    \caption{Increments}
%    \end{subfigure}
%    \caption{Cost and increments of one step of implicit Euler method using Algorithm \ref{alg:rkcd} as functions of $\Delta x$ and iteration number respectively along with theoretical bounds.}
%    \label{fig:euler}
%\end{figure}




\begin{figure}
    \centering
    \begin{subfigure}{0.49\linewidth}
    \centering
    \includegraphics[width=\textwidth]{cost1d.eps}
    \caption{1D case, DOF: $4,8,16,32,64$}
    \end{subfigure}
    \begin{subfigure}{0.49\linewidth}
    \centering
    \includegraphics[width=\textwidth]{cost3d.eps}
    \caption{3D case, DOF: $4^3,8^3,16^3,32^3,64^3$}
    \end{subfigure}
    \caption{Comparison of the cost (number of matrix-vector products)   as a function of the condition number $\kappa$ for Algorithm \ref{alg:3}. }
    %for computing  with prescribed tolerance $\varepsilon=10^{-10}$ one step of the implicit Euler method and SDIRK4 applied to  \eqref{eq:heat} for $N=1,3$ with meshsizes $\Delta t=5\times 10^{-2}$, $\Delta x= 1/4,1/8,1/16,1/32,1/64$, $\kappa=4 N\gamma\Delta t/\Delta x^2+1$,}
    \label{fig:3dLaplace}
\end{figure}
%In terms of simplicity, we discretize in space with a finite difference method with uniform grid size $\Delta x$. If we now denote with $y(t)=(y^{0}(t),\cdots  y^{N}(t)),  N\Delta x=1$, , then the semi-discrete formulation is given by the ODE
%\begin{equation} \label{eq:disc_form}
%\frac{dy}{dt}=-Dy, 
%\end{equation}
%where $-D$ is a symmetric matrix with real eigenvalues %$\lambda_{\text{heat}}$ in the interval 
%\[
% [-4a/(\Delta x)^{2},0 ]
%\]
%If we were to solve \eqref{eq:disc_form} with the implicit Euler method as we have seen in the previous section \eqref{eq:gf1} has a gradient structure with
%\[
%f(y)=\frac{\Delta t}{2}y^{T}Dy+ \frac12 \| y - y_n\|_2^2 
%\]
%Now  $f$ is quadratic with 
%\[
%C=I-\Delta t D, \quad  b=-y_{n} \quad c= \frac{1}{2}y^{T}_{n}y_{n}
%\]
%where the eigenvalues of $C$ are in $[\ell,L]$ where $\ell=1$ and $L=1+4\Delta t/(\Delta x)^{2}$. Now using Algorithm \ref{alg:rkcd}, Proposition \ref{prop:quad_RKC} guarantees the convergence of the iterates $x_{k}$ towards the minimizer $x_{*}=y_{n+1}$, the solution of the implicit Euler method at time $t_{n+1}$.
%In addition, given $y_{n}$ the cost  to approximate $y_{n+1}$   with target accuracy $\epsilon$ scales proportionally to $\mathcal{O}(\sqrt{\Delta t}/\Delta x \log{\epsilon^{-1}})$.
\paragraph{SDIRK4}
We  now consider the solution of \eqref{eq:main} by SDIRK4 when the different terms are given by \eqref{eq:pure_diffusion}. In this case $m=5$ and the Runge-Kutta matrix $\mathcal{A}$ is given by \eqref{eq:sdirk4}, hence $\gamma=1/4$ and
{
\begin{subequations} \label{eq:pure_diffusion_sdirk}
\begin{align} 
G_{D}(\mathbf{Y}) &= \mathbbm{1}_{m}\otimes y_{n}-\mathbf{Y}-\frac{1}{4}\Delta t (I_m \otimes D) \mathbf{Y}, \\ G_{A}(\mathbf{Y})&= \Delta t \left[\left(\mathcal{A}-\frac{1}{4} I_{m} \right)\otimes D \right]\mathbf{Y}   \\ G_{R}(\mathbf{Y})&=0.
\end{align}
\end{subequations}
}

Unlike the implicit Euler method, we see that in the case of SDIRK4 
Algorithm \ref{alg:explicitimplicit}  doesn't coincide with Algorithm 
\ref{alg:rkcd}, hence Proposition \eqref{prop:quad_RKC} does not 
directly apply. Nevertheless,  as we observe\footnote{we have used the 
same parameters $\varepsilon, \Delta x, \Delta t$ as for the implicit 
Euler} in  Figure \ref{fig:3dLaplace} given $y_{n}$, the cost to 
approximate $y_{n+1}$ with Algorithm \ref{alg:3} scales again 
proportionally to $\mathcal{O}(\sqrt{\kappa})$  independently of the 
value of $N$, where $\kappa$ is the condition number of $G_{D}$. We will prove in Theorem \ref{thm:diff} that this is indeed a property of Algorithm \ref{alg:3} for the implementation of an $m$-stage SDIRK method in the pure linear diffusion case, equation \eqref{eq:pure_diffusion}.




%Again here   $y(t)=(y^{0}(t),\cdots  y^{N}(t)),  N\Delta x=1$ with   $y^{i}(t) \approx u(x_{i},t)$ with $x_i=i\Delta x$.  In this case, $g(y(t))=-D y(t)$ and thus the corresponding ODE \eqref{eq:gf1} becomes
%\begin{equation} \label{eq:sdirk4_gf}
%\dot {\boldsymbol{Y}}(t)=-\boldsymbol{Y}(t)
%	+\mathbbm{1}_m \otimes y_n
%	-\Delta t (\calA \otimes D){\boldsymbol{Y}}(t)
%\end{equation}
%where ${\boldsymbol{Y}}(t) \in \mathbb{R}^{m \cdot N}$ is defined in \eqref{eq:summary}. Note that because the matrix $A$ \eqref{eq:sdirk4} of Runge-Kutta coefficients is non-symmetric then the right-hand side of \eqref{eq:sdirk4_gf} doesn't correspond to the gradient of a function, hence Algorithm \ref{alg:rkcd} is not directly applicable. On the other hand, the Jacobian of the right-hand side of \eqref{eq:sdirk4_gf} is given by 
%\begin{equation} \label{eq:naive_sdirk4}
%Q=-I_{m \cdot N}-\Delta t (\calA \otimes D).
%\end{equation}
%However, since the matrix $\calA$ \eqref{eq:sdirk4} is lower triangular with $1/4$ in the diagonal, the Hessian of the right-hand side of \eqref{eq:sdirk4_gf} has  negative real eigenvalues with multiplicity $m=5$ given by 
%\begin{equation}
%\lambda_{j}=-1-\frac{\Delta t}{4}\lambda_{\text{heat},j}     
%\end{equation}
%where $\lambda_{\text{heat},j}$, $j=1,\ldots,N$ are the eigenvalues of $D$. Hence one can consider using Algorithm \ref{alg:rkc} where $g$ is the right-hand side of \eqref{eq:sdirk4_gf} and where the time step used from the RKC integrator is now denoted by $h$. We choose the values of $s$ and $h$ according to \eqref{eq:defn of parameters} and \eqref{eq:defn of parameters1} where $\kappa$ is the condition number of the Hessian of the right-hand side of \eqref{eq:sdirk4_gf} with $\Delta t=5\times 10^{-2}$.  Figure \ref{fig:conv_sdirk4} shows that the increments (difference between iterates and exact solution of SDIRK4) do not converge to machine precision (blue curve).  We discuss this in depth in Section \ref{appendix:naive}, but in a nutshell, the issue is caused from the fact  that the matrix $Q$ in \eqref{eq:naive_sdirk4}  is non-symmetric.
        
 %       To deal with this issue, we consider splitting the coefficient matrix in the following way into a diagonal part and a non-diagonal part 
%\begin{subequations}        
%\begin{eqnarray}
%g_{\text{diag}}(\boldsymbol{Y}):&=&-(I_{md}+\gamma \Delta t (\mathbbm{1}_m \otimes D))\boldsymbol{Y}+\mathbbm{1}_m \otimes y_n \\
%g_{\text{non-diag}}(\boldsymbol{Y}):&=&-\Delta t ((\calA-\gamma \mathbbm{1}_m) \otimes D){\boldsymbol{Y}}.       
%\end{eqnarray}
%\end{subequations}
%Equation \eqref{eq:sdirk4_gf} becomes
%\begin{equation} \label{eq:split}
%\dot {\boldsymbol{Y}}(t)=g_{\text{diag}}(\boldsymbol{Y}(t))+g_{\text{non-diag}}(\boldsymbol{Y}(t)).
%\end{equation}
%Note that for nonlinear problems we have
%\begin{subequations}        
%\begin{eqnarray}
%g_{\text{diag}}(\boldsymbol{Y}):&=&-\boldsymbol{Y} + \gamma \Delta t g[\boldsymbol{Y}]+\mathbbm{1}_m \otimes y_n \\
%g_{\text{non-diag}}(\boldsymbol{Y}):&=&-\Delta t ((\calA-\gamma \mathbbm{1}_m) \otimes \mathbbm{1}_d){g[\boldsymbol{Y}]}.   
%\end{eqnarray}
%\end{subequations}
%Inspired by \cite[Algorithm 2]{EVVZ21}, we consider the following partitioned scheme (preserves steady state) to iterate on \eqref{eq:split}. Given the iterate $\boldsymbol{Y}_k$, apply the following recurrence to calculate $\boldsymbol{Y}_{k+1}$
%\begin{equation}
%	\begin{split}
%		K_0&=\boldsymbol{Y}_k,\\
%		K_1&=\boldsymbol{Y}_k+\mu_1h\left(g_{\text{diag}}(\boldsymbol{Y}_k)+g_{\text{non-diag}}(\boldsymbol{Y}_k)\right),\qquad \text{and for }\, j=2,\dots,s\\
%		K_j&=\mu_jh(g_{\text{diag}}(K_{j-1})+g_{\text{non-diag}}(\boldsymbol{Y}_k))+\nu_jK_{j-1}+(1-\nu_j)K_{j-2},\\
%		\boldsymbol{Y}_{k+1}&=K_s.
%	\end{split}
% \label{eq:partition1}
%\end{equation}
%The solution at time $t_{n+1}$ is estimated by $y_{n+1}=y_n+\Delta t \sum_{i=1}^s b_i g(Y_i)$. Note that in the case of SDIRK4 \eqref{eq:sdirk4} we have $y_{n+1}=Y_m$. 
%\begin{remark}
%    The recurrence \eqref{eq:partition1} is equivalent to standard RKC (see Algorithm \ref{alg:rkc}) applied to 
%\begin{equation} \label{eq:split1}
%\dot{\boldsymbol{Y}}(t)=g_{\text{diag}}(\boldsymbol{Y}(t))+g_{\text{non-diag}}(\boldsymbol{Y}_k).
%\end{equation}
%\end{remark}
%The recurrence \eqref{eq:partition1} is equivalent to standard RKC (see Algorithm \ref{alg:rkc}) applied to 
%\begin{equation} \label{eq:split1}
%\dot{\boldsymbol{Y}}(t)=g_{\text{diag}}(\boldsymbol{Y}(t))+g_{\text{non-diag}}(\boldsymbol{Y}_k)
%\end{equation}



%begin{figure}
%    \centering
%    \includegraphics[scale=0.3]{increment_split.png}
%    \caption{Increment versus iterations for the SDIRK4 method with and without partitioning.}
%    \label{fig:conv_sdirk4}
%\end{figure}
%The following theorem states that the partitioned RKC method applied with splitting to a linear ODE converges at a rate similar the quadratic case of Proposition \ref{prop:quad_RKC}. This is illustrated in Figure \ref{fig:conv_sdirk4} (black curve).

%{\color{red} \noindent For advection: $\|\boldsymbol{Y}_k-\boldsymbol{Y}_\star\|_2 \leq \kappa_2(PQ_\varepsilon)
%(\alpha_s(\eta) + \varepsilon \gamma^{-1}(2  + \eta/\ell)+\Delta t(\gamma+\varepsilon)\eta/\ell \|\Gamma\|_2)^k \|\boldsymbol{Y}_0-\boldsymbol{Y}_\star\|_2$.}

\subsubsection{Why the partitioning?}
We now discuss in detail why the partitioning \eqref{eq:part} is necessary for the construction of   Algorithm \ref{alg:3} and why applying  Algorithm \ref{alg:rkcd} without the gradient assumption to \eqref{eq:gff} would fail in general, even in the simplest possible case of pure linear diffusion where $F_{D}, F_{R}, F_{A}$ are given by \eqref{eq:pure_diffusion}.  More precisely, then  \eqref{eq:gff} becomes
\begin{equation}\label{eq:naive_sdirk4}
\dot{ \boldsymbol{Y}}(t)= \mathbbm{1}_{m} \otimes y_{n} +\mathcal{V} \boldsymbol{Y}(t), 
\quad 
\mathcal{V}\coloneqq-I_{md}-\Delta t (\calA \otimes D).
\end{equation}
Note that because $\calA$ is lower triangular with all entries in the diagonal equal to $\gamma$ when $m \geq 2$ neither $\calA$ nor $\mathcal{V}$ can be put 
in diagonal form. In addition, the eigenvalues of $\mathcal{V}$ all have multiplicity $m$ and are given by 
\[
\lambda_{j}=-1-\Delta t \gamma \lambda_{\text{heat},j}  
\]
where $\lambda_{\text{heat},j}$ are the eigenvalues of $D$, and the corresponding eigenspace associated to $\lambda_{j}$ has dimension $1$. 

We start by considering the case of Implicit Euler ($\gamma=1$, $m=1$)  for $\Delta t=5\times 10^{-2}, N=1$. Since now $G_{A}=G_{R}=0$ as discussed above Algorithm \ref{alg:3}  coincides with Algorithm \ref{alg:rkcd} and thus  Proposition \ref{prop:quad_RKC} applies. This is indeed the case up to round-off errors as we can see in Figure \ref{fig:iterations} where we plot the error between the iterations $x_{k}$ for Algorithm \ref{alg:3} and  the exact solution $\boldsymbol{Y}_{n+1}$ calculated using standard Newton iteration.  

If we now consider  the case of SDIRK4 ($\gamma=1/4$ and $m=5$), $G_{A} \neq 0$ and then \eqref{eq:gff} even though linear in $\boldsymbol{Y}$ is not of gradient form. As we can  see in  Figure \ref{fig:iterations}, and we  also prove in Theorem \ref{thm:diff}, the partitioning \eqref{eq:part} is essential for ensuring the convergence of iterates, since the iterates calculated with the non-partitioned Algorithm \ref{alg:rkcd} {(corresponding to Algorithm \ref{alg:3} with 
the definition of $G(\boldsymbol{x})$ replaced by $G(\boldsymbol{x})=G_D(\boldsymbol{x})+G_A(\boldsymbol{x})$ and with $G_R(\boldsymbol{x})=0$, or equivalently applying Algorithm \ref{alg:rkcd} to equation \eqref{eq:gff} with $F_A=F_R=0$)} suffer from   high-round off errors  that   deteriorate as stiffness increases {(see  Appendix \ref{appendix:naive} for details).} In contrast, iterates calculated with Algorithm \ref{alg:3} do not suffer from these issues. 
 \begin{figure}
    \centering
    \begin{subfigure}{0.48\linewidth}
    \centering
\includegraphics[width=\textwidth]{figconvexpSDIRK100.eps}
    \caption{$\Delta x= 10^{-2},\kappa=0.2\cdot 10^{4}$.}
    \end{subfigure}
    \begin{subfigure}{0.48\linewidth}
    \centering
    \includegraphics[width=\textwidth]{figconvexpSDIRK1000.eps}
    \caption{$\Delta x= 10^{-3},\kappa=0.2\cdot 10^{6}$.}
    \end{subfigure}
    \caption{Behaviour of the error  for Algorithm \ref{alg:3} and the non-partitioned method}
    
    %Convergence of the explicit stabilised iterations of Algorithm~\ref{alg:rkcd} for the implicit Euler method and the SDIRK4 method (with and without partioning) for computing one step with size $\Delta t=5\times 10^{-2}$ for the pure 1D diffusion problem \eqref{eq:heat}. We compare different condition number $\kappa=4 \gamma\Delta t/\Delta x^2+1$ with $\Delta =x=10^{-2},10^{-3}$.}
    \label{fig:iterations}
\end{figure}

We  explain this difference in depth in Section \ref{appendix:naive}, but in a nutshell, the issue is caused by the fact that the matrix $\mathcal{V}$ in \eqref{eq:naive_sdirk4}  is non-symmetric when $m \geq 2$  which creates issues for the non-partitioned method.  In the case of the partitioning \eqref{eq:part}, 
we obtain
\[
G_{D}(\boldsymbol{Y})=\mathbbm{1}_m \otimes y_n +\mathcal{V}_{1}\boldsymbol{Y},\qquad G_{A}(\boldsymbol{Y})=\mathcal{V}_{2}\boldsymbol{Y},\qquad
G_{R}(\boldsymbol{Y}) =0,
\]
where  the Jacobian $\mathcal{V}$ is decomposed in two parts $\mathcal{V}=\mathcal{V}_{1}+\mathcal{V}_{2}$ as
\begin{eqnarray*}
    \mathcal{V}_{1} &=& -I_{m}\otimes \hat{\mathcal{V}}_{1}, \quad \hat{\mathcal{V}_{1}}= (I_{d}+\gamma \Delta t D), \\
    \mathcal{V}_{2} &=&-\left(\mathcal{A}-\gamma I_{m} \right)\otimes \hat{\mathcal{V}}_{2}, \quad \hat{\mathcal{V}}_{2}=-\Delta t D.
\end{eqnarray*}
By construction, $\mathcal{V}_{1}$ is symmetric and has the same eigenvalues as $\mathcal{V}$, while $\mathcal{V}_{2}$ is non-symmetric and nilpotent. 

The iterates $\boldsymbol{x}_{k}$ for Algorithm \ref{alg:3} satisfy the following recurrence relation
\begin{eqnarray} 
    \boldsymbol{x}_{k+1}
   & = &
    R_s(h \mathcal{V}_{1})
    \boldsymbol{x}_k
   +
   hB_s(h \mathcal{V}_1)
   (\mathbbm{1}_m \otimes y_n +
	\mathcal{V}_2 \boldsymbol{x}_k) \label{eq:it_rkcd} \nonumber \\
&=&
    \boldsymbol{x}_k+ h B_s(h \mathcal{V}_{1})
   ( \mathbbm{1}_m \otimes y_n +
	\mathcal{V}\boldsymbol{x}_{k}) 
    %\quad \textcolor{red}{\mathbbm{1}_m \otimes y_n~SHOULD~BE~IN~G_A??}
    \label{eq:it_rkcd1}
\end{eqnarray}
where $B_s(z)=(R_s(z)-1)/z$, and we used $R_s(z)=1+B_s(z)z$. It is important to note that the steady state of \eqref{eq:it_rkcd1} coincides with the steady state of \eqref{eq:naive_sdirk4} as desired. Furthermore, as we can see in \eqref{eq:it_rkcd}, only the symmetric part of the matrix {$\hat{\mathcal{V}}$} is taken as an argument of the stability function $R_{s}$, while the non-diagonalizable part {$\hat{\mathcal{V}}_2$} is stabilised by $B_s$. This is, in fact, key for establishing the convergence properties of Algorithm \ref{alg:3} as detailed in the next section. 
\begin{remark}
Another benefit of the splitting $\mathcal{V}=\mathcal{V}_1+\mathcal{V}_2$ is that it allows for a parallel implementation of the stabilization terms $ R_s(h \mathcal{V}_{1})$ and $B_s(h \mathcal{V}_{1})$ over $m$ processors due to their block diagonal structures. Note that a similar splitting was proposed in \cite{VSC92} with the motivation of parallelizing the  iterates of a quasi-Newton implementation of SDIRK methods.
\end{remark}



%As we will see next, this decomposition of the Jacobian is essential for ensuring good convergence properties for Algorithm \ref{alg:explicitimplicit} in the case of pure linear diffusion for a general $m$-stage DIRK method 



 



\subsubsection{Theoretical analysis of Algorithm \ref{alg:explicitimplicit}} \label{subsubsec:analysis}

\begin{theorem}\label{thm:diff}
Consider Algorithm \ref{alg:explicitimplicit} for the implementation of an $m$-stage SDIRK  method applied to ODE \eqref{eq:main}  where  $F_{D}, F_{A}, F_{R}$ are given by \eqref{eq:pure_diffusion}.
Then, for all $\varepsilon>0$ there exists $C_\varepsilon$ (independent of the dimension $d$) such that the iterates $\boldsymbol{x}_k$ of Algorithm \ref{alg:explicitimplicit} converge to the
SDIRK solution $\boldsymbol{Y}_{n+1}$ of \eqref{eq:nonlinearimp} with the following convergence estimate,
\begin{equation} \label{eq:proof1}
\|
\boldsymbol{x}_k - \boldsymbol{Y}_{n+1}
\|_2 \leq C_\varepsilon \big(\alpha_{s}(\eta) + \varepsilon\big)^k 
\|
\boldsymbol{x}_0 - \boldsymbol{Y}_{n+1}
\|_2.
\end{equation}
\end{theorem}
\begin{proof}
%We recall that for  SDIRK, the Runge-Kutta matrix of coefficients $\calA$ is lower triangular with coefficients $a_{ii}=\gamma$ all equal on the diagonal.
%Applying the partitioned Chebyshev method with time step $h$ to the system of ODE
%\begin{equation}\label{eq:gf2}
%			\dot{\boldsymbol{Y}}=
%			(-\Delta t (\gamma \mathbbm{1}_m \otimes D) -I_{md})
%%			+\left(-\Delta t((\calA-\gamma \mathbbm{1}_m) \otimes D)
%			\boldsymbol{Y}
%			+\mathbbm{1}_m\otimes y_n\right).%		\end{equation}
%where $\boldsymbol{Y}(0)=\boldsymbol{Y}_0$ is the initial condition, and $\boldsymbol{Y}_\star$ is the steady state,
%$$
%0 = (-\Delta t (\gamma \mathbbm{1}_m \otimes D) -I_{md})
%			\boldsymbol{Y}_\star
%			+\left(-\Delta t((\calA-\gamma \mathbbm{1}_m) \otimes D)
%			\boldsymbol{Y}_\star
%			+\mathbbm{1}_m\otimes y_n\right)
%$$
Using \eqref{eq:it_rkcd}, we obtain,
%We have an induction of the form
%\begin{equation}
%\boldsymbol{x}_{k+1} - \boldsymbol{Y}_{n+1}= 
%\begin{pmatrix}
%				R(h Q_{1}) & 0 & \ldots & \ldots & 0\\
%				a_{21}B(hQ_{1})hQ_{2} & R(z_1) & \ldots & \ldots & 0\\
%				a_{31}B(z_1)z_2 & a_{32}B(z_1)z_2 & R(z_1) & \ldots & 0\\
%				\vdots & \vdots &&\ddots & \vdots\\
%				a_{m1}B(z_1)z_2 & a_{m2}B(z_1)z_2 & \ldots & a_{m-1,m}B(z_1)z_2 & R(z_1)
%			\end{pmatrix}
% (\boldsymbol{x}_{k}- \boldsymbol{Y}_{n+1})
%\end{equation}
%where $z_1=-h(\gamma \Delta t D + \mathbbm{1}_d)$, $z_2=-h\Delta t D$,  $B(z)=(R(z)-1)/z$, or equivalently
\begin{equation}
\boldsymbol{x}_{k+1} - \boldsymbol{Y}_{n+1}= M (\boldsymbol{x}_{k}- \boldsymbol{Y}_{n+1}),\qquad M=  R_{s}(h\mathcal{V}_1) + B_{s}(h\mathcal{V}_1)h\mathcal{V}_2,
\end{equation}
where   $B_{s}(z)=(R_{s}(z)-1)/z$. In order to prove \eqref{eq:proof1} we need to have an estimate of $\norm{M}$. The main difficulty in doing so is the term $\mathcal{V}_{2}$, as this is not diagonalizable.  To address this issue, we construct an appropriate change of coordinates which makes this term negligible. Precisely, we consider the Jordan decomposition of the Runge-Kutta matrix $\calA$ of coefficients,
$
\calA = P\mathcal{J}P^{-1}
$
where $P\in\R^{m\times m}$ is an invertible matrix, and $\mathcal{J}\in\R^{m\times m}$ is a Jordan block matrix of the form
\begin{equation}\label{eq:defJ}
\mathcal{J}= 
			\begin{pmatrix}
				\gamma & 1 & 0& \cdots & 0\\
				0& \gamma & 1& \cdots & 0\\
				\vdots&\vdots&\ddots&\\
				0& 0 & & \ddots & 1 \\
				0& 0 &  &\cdots & \gamma \\
			\end{pmatrix} =\gamma I_m + F. 
\end{equation}
Here, we assume for simplicity of the presentation that there is a single Jordan block $\mathcal{J}$, but the proof generalises straightforwardly to the case of multiple Jordan blocks.

Now let $Q_\varepsilon=\mathrm{diag}(1,\varepsilon,\varepsilon^2,\ldots, \varepsilon^{m-1})$.
Using 
		$
		Q_\varepsilon^{-1} F Q_\varepsilon = \varepsilon F
		$,
we deduce
$
Q_\varepsilon^{-1}P^{-1} \calA PQ_\varepsilon =  \gamma \mathbbm{1}_m +  \varepsilon F
$
and
$$
(Q_\varepsilon^{-1}P^{-1}\otimes I_d) M (PQ_\varepsilon\otimes I_d) =  R_s(h\mathcal{V}_1) - \varepsilon   B_s(h\mathcal{V}_1) [F \otimes h\hat{\mathcal{V}}_{2}],
$$
where we have used that $(A \otimes B)(C \otimes D)=AC \otimes BD$ and that the matrix $PQ_\varepsilon\otimes I_d$ and its inverse both commute with $R(h\mathcal{V}_1)=I_m \otimes R_s(h\hat{\mathcal{V}}_{1})$ and $B_s(\mathcal{V}_1)$. This yields
\begin{eqnarray}
\|\boldsymbol{x}_k-\boldsymbol{Y}_{n+1}\|_2 &=& \|M^k(\boldsymbol{x}_0-\boldsymbol{Y}_{n+1})\|_2 \nonumber \\
&\leq& \kappa_2(PQ_\varepsilon)
(\|R(h\hat{\mathcal{V}}_{1})\|_2 + \varepsilon \|F\|_2 \|B_{s}(h\hat{\mathcal{V}}_{1})h\hat{\mathcal{V}}_{2}\|_2)^k \|\boldsymbol{x}_0-\boldsymbol{Y}_{n+1}\|_2
\end{eqnarray}
where $\kappa_2(PQ_\varepsilon)=\|PQ_\varepsilon\|_2\|(PQ_\varepsilon)^{-1}\|_2$ is the condition number of the matrix $PQ_\varepsilon\in\R^{m\times m}$, which is independent of $d$.
{ We now have (see \cite[Lemma A.1]{EVVZ21})
\begin{equation} \label{eq:RBone}
\|R_{s}(h\hat{\mathcal{V}}_1)\|_2 \leq \alpha_s(\eta)<1,\qquad \|B_{s}(h\hat{\mathcal{V}}_1)\|_2 \leq 1, 
\end{equation}
In addition, we have   $\|F\|_2 = 1$, 
 $\hat{\mathcal{V}}_2=\gamma^{-1}(I_d -\hat{\mathcal{V}}_1)$,
and $h\leq\eta e^\eta/\ell$ using \eqref{eq:h_bound},
which yields $\|B(h\hat{\mathcal{V}}_1)h\hat{\mathcal{V}}_1\|_2 = \|R_s(h\hat{\mathcal{V}}_1) -I_d\|_2 \leq 2$, and 
\begin{equation} \label{eq:us_bound}
\|B_s(h\hat{\mathcal{V}}_1)h \hat{\mathcal{V}}_2\|_2 \leq  \|h \hat{\mathcal{V}}_2\|_2 \leq\gamma^{-1}(2  + \eta e^\eta/\ell )
\end{equation}
}
and hence
\begin{eqnarray*}
\|\boldsymbol{x}_k-\boldsymbol{Y}_{n+1}\|_2 &\leq& \kappa_2(PQ_\varepsilon)
(\alpha_s(\eta) + \varepsilon \gamma^{-1}(2  + \eta e^\eta/\ell ))^k \|\boldsymbol{x}_0-\boldsymbol{Y}_{n+1}\|_2 \\
&\leq&C_\varepsilon
(\alpha_s(\eta) + \varepsilon)^k \|\boldsymbol{x}_0-\boldsymbol{Y}_{n+1}\|_2 
\end{eqnarray*}
where $C_\varepsilon = \kappa_2(PQ_{\varepsilon'})$,
with
$\varepsilon' = \varepsilon\gamma(2  + \eta e^\eta/\ell )^{-1} $ is a constant independent of the dimension $d$
and this concludes the proof.
\end{proof}
{In the proof of Theorem \ref{thm:diff}, observe that one can improve the estimate \eqref{eq:us_bound} for the quantity $B(h\hat{\mathcal{V}_1})h \hat{\mathcal{V}}_2$}, but this is not useful here because $\varepsilon$ can be chosen arbitrarily small. 
\begin{remark}
In the case of the implicit Euler method as we discussed before Algorithm \ref{alg:3} is equivalent to Algorithm \ref{alg:rkcd}. This is reflected in the statement of Theorem \ref{thm:diff} since $\varepsilon$=0, $C_{\varepsilon}=1$ in \eqref{eq:proof1} which agrees with \eqref{eq:ratte}  in Proposition \ref{prop:quad_RKC}.    
\end{remark}





\subsection{The case of linear  advection-diffusion}

We now consider  the 
following  linear advection-diffusion equation
\begin{equation} \label{eq:advdiffexact}
\partial_t u=a\Delta u+b \cdot \nabla u, \quad (t,x)\in [0,+\infty)\times (0,1)^{N}
\end{equation}
with periodic boundary conditions and a smooth initial condition $u(x,0)$, where $a>0$ and $b\in \R^N_+$ are fixed.  We  now  discretise in space using a finite difference method with a uniform mesh size $\Delta x$ in each spatial dimension.  We thus have 
\begin{equation} \label{eq:adv_diff}
F_{D}(y)= -Dy, \quad F_{A}(y)=E y \quad F_{R}(y)=0.
\end{equation}
where $D$ is the discrete Laplacian with periodic boundary conditions and $E$ is the discrete centred advection operator.
Following the presentation in \cite{Zb11}  the eigenvalues $\lambda_{\text{heat-adv},j}=\lambda_{\text{heat},j}+i\mu_{\text{adv},j}$ of $F=F_D+F_A$ are located  on the following ellipse in the complex plane
\begin{equation} \label{eq:ellipse_pq}
\left(\frac{2p}\alpha+1\right)^2 + \left(\frac{q}\beta\right)^2=1,
\end{equation}
with half-height $\beta=\norm{b}_{1}\Delta x^{-1}$ and width $\alpha=4aN\Delta x^{-2}$, and where
$p=\lambda_{\text{heat},j}\in[-\alpha,0]$ located on the negative real axis are the eigenvalues of $F_D$ and $iq=i\mu_{\text{adv},j}\in[-i\beta,i\beta]$ on the imaginary axis are the eigenvalues of $F_A$.

Before we present a theorem for an $m$-stage SDIRK method it is useful to introduce the following function
\begin{equation}\label{eq:stabRIE}
\mathcal{R}_{s}(P,Q) =  R_s(P) + B_s(P)Q.
\end{equation}
that will play an important role in establishing the stability and convergence properties in the general case. For simplicity, we start by considering the implicit Euler case for which the iteration error for Algorithm \ref{alg:3} satisfies 
%to detail first the case of the implicit Euler. In this case, have
%\[
%G_{D}(y)=y_{n}+Q_{1}y, \quad G_{A}(y)= Ey, \quad G_{R}(y)=0 
%\]
%We then obtain the partitioning
%\[
%G_{D}(\boldsymbol{Y})= \mathbbm{1}_{m}\otimes y_{n}-\mathbf{Y}-\gamma\Delta t [\mathbbm{1}_m \otimes D] \mathbf{Y}, \quad G_{A}(\boldsymbol{Y})=[{\calA\otimes}E]\boldsymbol{Y}+\Delta t \left[\left(\mathcal{A}-\gamma I_{m} \right)\otimes D \right]\boldsymbol{Y}, \quad G_{R}(\boldsymbol{Y})=0
%\]
%Before we proceed with providing a theorem for the general case of an $m$-stage SDIRK method, we provide a stability analysis for the case of the implicit Euler method. More precisely, in this case, we have 
\begin{equation}\label{eq:iterrIE}
x_{k+1}-y_{n+1} = \mathcal{R}_{s}(-hI_d - h\Delta t D,hE)(x_k-y_{n+1})
\end{equation}
and consider the related stability domain 
\begin{equation} \label{eq:stab_domain}
{\mathcal{S}_\gamma} = \{(p,q) \in \mathbb{R}^2\ ;\ R_{\text{stab}}({\gamma} p,q) <1 \}.
\end{equation}
{with $\gamma=1$ for the implicit Euler method and} where
\[
R_{\text{stab}}(p,q)\coloneqq |R_{s}(-h+h p)|^2 + |B_{s}(-h+h p)|^2|h q|^2
\]
We first prove the stability property of Algorithm \ref{alg:3} in the simplest case of the implicit Euler method and based on the stability domain \eqref{eq:stab_domain}. We then show that it also implies the convergence property of Algorithm \ref{alg:3} for an $m$-stage SDIRK method. 

\begin{proposition}\label{prop:rate1}
The error $\norm{x_{k+1}-y_{n+1}}_{2}$ for the implicit Euler method given by \eqref{eq:iterrIE} converges to zero as $k\rightarrow+\infty$ if for all $p_{j}=-\Delta t \lambda_{\text{heat},j}, \ q_{j}=\Delta t\mu_{adv,j}, j=1,\ldots,d$, we have $(p_{j},q_{j}) \in \mathcal{S}_1$. In this case, we have
$$
\norm{x_{k+1}-y_{n+1}}_{2} \leq \rho\norm{x_{k}-y_{n+1}}_{2},\qquad \rho := \max_{j=1,\cdots,d} R_{\text{stab}}(p_j,q_j)^{1/2}<1.
$$
\end{proposition}
\begin{proof}
If $p_{j},q_{j} \in \mathcal{S}_1$ for all $j$ we have that
 $R_{\text{stab}}(p_{j},q_{j})<1$. Now, since we assume periodic boundary conditions, the matrices   $D,E$ commute which implies that  
$ 
 \|\mathcal{R}_{s}(-hI_d - h\Delta t D,hE) \|_2 = \rho
 <1,
$
which yields the convergence for the error $\norm{x_{n+1}-y_{n+1}}_{2}$ in \eqref{eq:iterrIE}.     
\end{proof}





In proving Proposition \ref{prop:rate1}, we have explicitly taken into account that because of the periodic boundary conditions, the matrices $D$ and  $E$ are commuting. This will, in general, not hold, for example in the case of Neumann or Dirichlet boundary conditions. It is still possible to establish a stability criterion, as stated in the next proposition. 

\begin{proposition} \label{prop:rate2}
Consider \eqref{eq:main} with $F_D,F_A,F_R$ given  in \eqref{eq:adv_diff} and where $D$ is a given positive semi-definite matrix and $E$ is skew symmetric matrix with bound $\norm{E} \leq \norm{b}_{1}\Delta x^{-1}$. Then the  error $\norm{x_{k+1}-y_{n+1}}_{2}$ for Algorithm \ref{alg:3} applied with the implicit Euler method given by \eqref{eq:iterrIE} converges to zero as $k\rightarrow+\infty$ if 
 \begin{equation} \label{eq:rate_estimate}
 \hat \rho := \alpha_s(\eta) + \frac{\eta e^\eta}{\ell}\norm{b}_{1}\frac{\Delta t}{\Delta x} < 1.
 \end{equation}
In this case we have
\[
\norm{x_{k+1}-y_{n+1}}_{2} \leq \hat \rho\norm{x_{k}-y_{n+1}}_{2}.
\]
\end{proposition}
\begin{proof} 
Using \eqref{eq:RBone}
%$|R_{s}(-h+h p)|< \alpha_s(\eta)$, $|B_{s}(-h+h p)|\leq 1$, 
and $\norm{E} \leq \norm{b}_{1}\Delta x^{-1}$, 
we obtain from \eqref{eq:iterrIE}, $ \|\mathcal{R}_{s}(-hI_d - h\Delta t D,hE) \|_2\leq  \alpha_s(\eta) + h \Delta t \norm{E} \leq \hat  \rho <1$ where we use  \eqref{eq:h_bound},\eqref{eq:rate_estimate} and the convergence holds.
\end{proof}




 





We now state the general result, that the estimate for implicit Euler generalises for an arbitrary $m$-diagonal SDIRK method. 




%In terms of simplicity, we discretize space with a finite difference method with uniform grid size $\Delta x$. If we now denote with $y(t)=(y^{0}(t),\cdots  y^{N}(t)),  N\Delta x=1$, where  $y^{i}(t) \approx u(x_{i},t)$ with $x_i=i\Delta x$, then the semi-discrete formulation is given by the ODE
%\begin{equation} \label{eq:disc_form_ad}
%\frac{dy}{dt}=-Dy+\Gamma y, 
%\end{equation}
%where $D$ is a symmetric positive-definite matrix whose eigenvalues grow as $1/\Delta x^2$ and $\Gamma$ is a skew-symmetric matrix whose eigenvalues grow as $1/\Delta x$ on the imaginary axis. Taking advantage of the damping provided by the polynomial $B(z)$, the implementation of the method for advection-diffusion problems is done by applying Algorithm \ref{alg:expSDIRK} where the advection term is treated in the same way as $g_{\text{non-diag}}$. For linear problems, we have the following theorem.

\begin{theorem}\label{thm:ad}
Consider Algorithm \ref{alg:3} for the implementation of an $m$-stage SDIRK Runge-Kutta method applied to ODE \eqref{eq:main} where $F_{D},F_{A},F_{R}$ are given by \eqref{eq:adv_diff}. Assume that $D$ and $E$ satisfy either the assumptions of Proposition \ref{prop:rate2} or the stronger assumptions of Proposition \ref{prop:rate1}, and
{assume that at least one of the following bounds hold}
{\begin{align*}
\hat \rho:&= \alpha_s(\eta)+ \frac{\eta e^\eta}{\ell}\norm{b}_{1}\frac{\Delta t}{\Delta x}<1,
 \\
\rho:&= \max_{j=1,\ldots,d} \left(\mathcal{R}_s(-h-h\gamma\Delta t \lambda_{\text{heat},j},h\Delta t \mu_{\text{adv},j})^{1/2}\right)<1.
\end{align*}}
Then, for all $\varepsilon>0$ there exists $C_\varepsilon$ (independent of the dimension $d$)  such that the iterates $\boldsymbol{x}_k$ of Algorithm \ref{alg:explicitimplicit} converge to the
SDIRK solution $\boldsymbol{Y}_{n+1}$ of \eqref{eq:nonlinearimp} with the following convergence estimate,
    \begin{equation} \label{eq:boundadv}
\|\boldsymbol{x}_k-\boldsymbol{Y}_{n+1}\|_2 \leq C_\varepsilon
\left(\varepsilon+ \min(\rho,\hat \rho)
\right)^k \|\boldsymbol{x}_0-\boldsymbol{Y}_{n+1}\|_2.
\end{equation} 
%$$
%\|\boldsymbol{x}_k-\boldsymbol{Y}_{n+1}\|_2 \leq C_\varepsilon\left(\alpha_s(\eta) + \epsilon+ \frac{\gamma\eta}{\ell}b\frac{\Delta t}{\Delta x}\right)^k \|\boldsymbol{x}_0-\boldsymbol{Y}_{n+1}\|_2.$$
\end{theorem}
\begin{proof}
    Following the lines of the proof of Theorem \ref{thm:diff}, we obtain
$$
\boldsymbol{x}_{k+1} - \boldsymbol{Y}_{n+1}= (M+\tilde M) (\boldsymbol{x}_{k}- \boldsymbol{Y}_{n+1}).
$$
%\quad M=  R_{s}(h\mathcal{V}_1) + B_{s}(h\mathcal{V}_1)(h\mathcal{V}_2),\ \tilde M=B_{s}(h\mathcal{V}_1)(\mathcal{A}\otimes h\Delta t E),
Here $M=R_{s}(h\mathcal{V}_1) + B_{s}(h\mathcal{V}_1)(h\mathcal{V}_2)$ is exactly the same matrix that appears in Theorem \ref{thm:diff} while the new term $\tilde M=B_{s}(h\mathcal{V}_1)(\mathcal{A}\otimes h\Delta t E)$ arises due to the presence of advection.
Now using the same change of coordinates as in Theorem \ref{thm:diff} we have that 
\begin{align*}
(Q_\varepsilon^{-1}P^{-1}\otimes I_d) (M+\tilde M) (PQ_\varepsilon\otimes I_d) &= I_m \otimes R(h\hat{\mathcal{V}}_1) 
+ \varepsilon F \otimes B(h\hat{\mathcal{V}}_1)h\hat{\mathcal{V}}_2 \\
&+ (\gamma I_m+\varepsilon F)\otimes B(h\hat{\mathcal{V}}_1)h\Delta t E,
\end{align*}
which yields
    \begin{align*} 
\|\boldsymbol{x}_k-\boldsymbol{Y}_{n+1}\|_2 &\leq C_\varepsilon
\Big(
%\alpha_s(\eta) +\epsilon \|F\|_2 \|B_s(h\hat{\mathcal{V}}_1)h\hat{\mathcal{V}}_2\|_2
 \|\mathcal{R}_{s}(-hI_d - h\Delta t \gamma D,hE) \|_2 \\
&+\varepsilon\|F\|_2 (\|B_s(h\hat{\mathcal{V}}_1)h\hat{\mathcal{V}}_2\|_2 + h\Delta t\|B_s(h\hat{\mathcal{V}}_1)\|_2\|E\|_2)
\Big)^k \|\boldsymbol{x}_0-\boldsymbol{Y}_{n+1}\|_2.
%& \leq C_\varepsilon
%\left(\alpha_s(\eta) + \varepsilon+\left(\frac{\gamma\eta e^\eta}{\ell}\right) %b\frac{\Delta t}{\Delta x}\right)^k  \|\boldsymbol{x}_0-\boldsymbol{Y}_{n+1}\|_2.
\end{align*} 
Now proceeding in a similar way as the end of the proof of Theorem \ref{thm:diff},
 we obtain the desired estimate \eqref{eq:boundadv}.
%     We obtain a similar induction for the iterations of the vector of internal stages of the SDIRK method, where the amplification matrix is now defined by
%     $$M= \mathbbm{1}_m \otimes R(z_1) + (\calA-\gamma \mathbbm{1}_m) \otimes B(z_1)z_2+\calA\otimes B(z_1)q$$ where $z_1$ and $z_2$ are defined as before and $q=h\Delta t\Gamma$. Hence we have
%     $$(Q_\varepsilon^{-1}P^{-1}\otimes \mathbbm{1}_d) M (PQ_\varepsilon\otimes \mathbbm{1}_d) = \mathbbm{1}_m \otimes R(z_1) + \varepsilon F \otimes B(z_1)z_2 + (\gamma \mathbbm{1}_m+\varepsilon F)\otimes B(z_1)q,
% $$
% which yields 
% \begin{align}
%     \|\boldsymbol{Y}_k-\boldsymbol{Y}_\star\|_2 &\leq \kappa_2(PQ_\varepsilon)
% (\|R(z_1)\|_2 + \varepsilon \|F\|_2\|B(z_1)z_2\|_2+(\gamma+\varepsilon\|F\|_2)\|B(z_1)\|_2 \|q\|_2)^k \|\boldsymbol{Y}_0-\boldsymbol{Y}_\star\|_2 \nonumber\\
%     &\leq \kappa_2(PQ_\varepsilon)
% \left(\alpha_s(\eta) + \varepsilon \gamma^{-1}(2  + \eta/\ell)+\frac{\eta}{\ell} (\gamma+\varepsilon)\Delta t\frac{b}{\Delta x}\right)^k \|\boldsymbol{Y}_0-\boldsymbol{Y}_\star\|_2 \nonumber\\
% &\leq C_\varepsilon
% \left(\alpha_s(\eta) + \epsilon+\left(\frac{\gamma\eta}{\ell}\right) b\frac{\Delta t}{\Delta x}\right)^k \|\boldsymbol{Y}_0-\boldsymbol{Y}_\star\|_2.
% \end{align}
%The last estimate is obtained in a similar way as the estimate of Theorem \ref{thm:diff}, where here we use in addition that $\|E\|_2\leq b/\Delta x$ and we define $C_\varepsilon = \kappa_2(PQ_{\varepsilon'})$ with $\varepsilon'=\varepsilon (\gamma^{-1} (2+\eta e^\eta/\ell)+\eta\Delta t b/\ell\Delta x)^{-1}$.
\end{proof}
Note that in the absence of advection ($b=0$), the estimate of Theorem \ref{thm:ad} reduces to the one of Theorem \ref{thm:diff} with $\hat \rho \leq \rho=\alpha_s(\eta)$. 

%\begin{remark}
 %   The bound in Theorem \ref{thm:ad} shows that for advection-diffusion problems, the iteration converges for $\Delta t$ small enough. As we will illustrate in the next section, this bound is not sharp and the method performs much better even when $\Delta t$ is quite large {\color{red}(...)}. 
%\end{remark}



\begin{figure}[t]
    \centering
    \begin{subfigure}{0.45\linewidth}
        \includegraphics[width=\linewidth]{stabdom.png}
        \caption{Stability domain}
    \end{subfigure}
    \begin{subfigure}{0.46\linewidth}
            \includegraphics[width=\linewidth]{zoomed_stabdom.png}
        \caption{Zoomed version}
    \end{subfigure}
    \caption{{Stability regions $\mathcal{S}_\gamma$ in \eqref{eq:stab_domain} with $\gamma=1/4$ of the diffusion-advection coupling for $s=20$ and damping parameters $\eta=4$. We also include the largest ellipse of the form \eqref{eq:ellipse_pq} tangent to the $y$-axis that can be included in the stability region $\mathcal{S}_\gamma$ in \eqref{eq:stab_domain} (blue dashed lines) where the iteration error decays, and in the stability region $\widehat{\mathcal{S}}_\gamma$ in \eqref{eq:stabdomainmod} where the iteration error is halved at each iteration.}}
    \label{fig:stability1}
\end{figure}
We now plot in Figure \ref{fig:stability1} the stability domain $\mathcal{S}_\gamma$ given by \eqref{eq:stab_domain} with $\gamma=1/4$ for $s=20$ and $\eta=4$. This domain covers a large  portion of the negative real axis as expected from the properties of the Chebyshev polynomials. On the other hand, the width on the $q$-axis near the origin is much narrower.  
{The blue ellipse in Figure \ref{fig:stability1} indicates the largest ellipse of the form \eqref{eq:ellipse_pq} that contains the origin with principal axes parallel to the $p$-$q$ axes and that can be included in the stability domain $\mathcal{S}_\gamma$.  If we denote by $d_s$ and $a_s$  respectively the width and the half-height of this ellipse  we observe that $d_{s}$ grows quadratically with the number of stages $s$ while  $a_{s}$ grows linearly with $s$ as illustrated in Figure \ref{fig:dsas}. More precisely, we find numerically for this range of parameters that $d_s$ and $a_s$ grow like 
\begin{equation} \label{eq:ellipsegrowth}
d_s = 2s^2,\qquad a_s = 0.62s.
\end{equation}
For comparison, we also plot in  in Figure \ref{fig:stability1} in red dotted lines the largest ellipse of the form \eqref{eq:ellipse_pq} that is included in the modified stability domain 
\begin{equation} \label{eq:stabdomainmod}
\widehat{\mathcal{S}}_\gamma = \{(p,q) \in \mathbb{R}^2\ ;\ R_{\text{stab}}(p,q) <1/2 \}
\end{equation}
where the iteration error is at least halved at each iteration using $\hat \rho<1/2$ in \eqref{eq:boundadv}. This is different to blue ellipse corresponding to the stability domain $\mathcal{S}$ in \eqref{eq:stab_domain} where for $\hat \rho<1$ only the decay of the iteration error is guaranteed. This illustrates the favorable stabilization of the explicit stabilised iteration taking advantage of the diffusion term even for advection terms with relatively large Peclet numbers.
}  
{Additionally, note that in the case where ${\rho}<1$ in Theorem \ref{thm:ad} is an appropriate bound for an SDIRK method, the behaviour of the stability function of the implicit Euler studied in Figures \ref{fig:stability1}-\ref{fig:dsas} dictates the choice of parameters, with  the only difference is that} {the stability domain size is scaled with respect to the diffusion {(horizontal $p$-axis in Figure \ref{fig:stability1})} by a factor $\gamma^{-1}$} %now the physical time is divided by $\gamma^{-1}$.
\begin{remark}
{As emphasized in \cite[Remark 3.4]{AV13} for the stability analysis of the PIROCK method for advection-diffusion problems, the linear growth in \eqref{eq:ellipsegrowth} of the ellipse half-height is a feature of the proposed explicit stabilised implementation of SDIRK methods as it allows for relatively large Peclet numbers.
Indeed, considering the advection-diffusion problem \eqref{eq:advdiffexact} with spatial discretization \eqref{eq:adv_diff}.
Substituting $h\Delta t \alpha = 2 s^2$ into $h\Delta t \beta \leq \max(0.62s,1)$  with $h\simeq 1.40$ yields 
$\Delta t \beta \leq \max(\frac{0.62^2}{2\cdot 1.40} \alpha/\beta, 1)$. 
Using $\beta=\norm{b}_{1}\Delta x^{-1}$, $\alpha=4aN\Delta x^{-2}$ we deduce the the stability sufficient condition%
$$\Delta t \leq \max(0.55 aN/\norm{b}_{1}, \Delta x/\norm{b}_{1}),$$
which is a stability constraint on the stepsize $\Delta t$ independent of
the spatial grid size  small enough ($\Delta x \leq 0.55 aN$). This is a feature of the proposed partitioned explicit stabilised implementation, shared with the PIROCK method satisfying an analogous stability sufficient condition (see \cite[Remark 3.4]{AV13}), in contrast to the other stabilised explicit integrators RKC, PRKC and ROCK2 where only small Peclet numbers can be considered.
}
\end{remark}

\begin{figure}[t]
    \centering
    \begin{subfigure}{0.43\linewidth}
        \includegraphics[width=\linewidth, height=4.3cm]{ds_eta4.png}
        \caption{Major axis $d_s$}
    \end{subfigure}
    \begin{subfigure}{0.43\linewidth}
        \includegraphics[width=\linewidth, height=4.3cm]{as_eta4.png}
        \caption{Minor radius $a_s$}
    \end{subfigure}
   
    \caption{{
    Stability of the diffusion-advection coupling for $\gamma=1/4,s=20$ and damping parameters $\eta=4$.
    Growth with respect to the stage parameter $s$ of the width $d_s$ and the half-height $a_s$ in \eqref{eq:ellipsegrowth} of the largest ellipse \eqref{eq:ellipse_pq} tangent to the y-axis that can be included in the stability region $\mathcal{S}_\gamma$ in \eqref{eq:stab_domain} with $\gamma=1/4$.}
    }
    \label{fig:dsas}
\end{figure}





%\subsection{Advection--diffusion--reaction equations {\color{red} TO BE CHECKED}}
% We now consider  the one-dimensional linear advection-diffusion equation
% \begin{equation}
% \partial_t u=a\Delta u+b \nabla u, \quad (t,x)\in[0,+\infty]\times[0,1]
% \end{equation}
% with periodic boundary conditions and a smooth initial condition $u(x,0)$. In terms of simplicity, we discretize space with a finite difference method with uniform grid size $\Delta x$. If we now denote with $y(t)=(y^{0}(t),\cdots  y^{N}(t)),  N\Delta x=1$, where  $y^{i}(t) \approx u(x_{i},t)$ with $x_i=i\Delta x$, then the semi-discrete formulation is given by the ODE
% \begin{equation} \label{eq:disc_form}
% \frac{dy}{dt}=-Dy+\Gamma y, 
% \end{equation}


% \begin{itemize}
% \item Present the method (31) in the case of advection-diffusion 
% \item Have the proof that contains $(\alpha_{s}(\eta)+\epsilon+\Delta t ||\Gamma||)^{k}$
% \item Explain how the one-dimensional test problem can give you the value of limiting $b$ and show case that we can do better because of the presence of $B(z)$
% \end{itemize}

% \[
% b \frac{\Delta t}{\Delta x} < 1-\alpha_{s}(\eta)
% \]

% \[
% b \frac{\Delta t}{\Delta x}<0.6 s
% \]


% The current approach involves applying an explicit stabilised integrator to equation \eqref{eq:gf1} using a time step $h=\bigo(1)$ and iterating through (Algorithm \ref{algo:diff}). 
  
%   For purely diffusive problems, competing with existing stabilised schemes presents challenges due to the need for iteration to calculate the solution at each step. As a result, the total cost will be the number of stages of the first-order RKC method multiplied by the number of iterations. Typically, our objective is to develop an integrator of order greater than two for problems with highly stiff reaction terms. This pursuit is driven by the fact that the best explicit stabilised methods available in the literature for handling stiff reactions are generally limited to order two at most \cite{AV13}.
  
  % In \cite[Section\ts4]{EVVZ21}, the authors introduce a modified algorithm designed to address minimization problems featuring a strongly convex objective function perturbed by a nonlinear, non-stiff, but potentially costly smooth convex function. 
 % We want to extend this approach to tackle problems encompassing diffusion, advection, and highly stiff reaction terms. An implicit treatment is necessary for the reaction term. This implicit treatment is applied exclusively to the final gradient flow and integrated with the first-order RKC method within a partitioning framework. Specifically, we utilize an implicit Euler method paired with a quasi-Newton iteration. Despite concerns of computational expense, this approach remains economical, as the reaction terms do not involve any derivatives of the unknown function. Consequently, their Jacobian matrix is block-diagonal, facilitating straightforward inversion.
  
 % Ideally, we would like to solve an ODE of the form
%		\begin{equation}
%			\dot y =F_D(y)+F_A(y)+F_R(y),
%			\label{eq:adrodr}
%		\end{equation}
%	where the vector fields $F_D$, $F_A$, and $F_R$ are obtained respectively from the discretization of a diffusion term, an advection term, and a reaction term in a time-dependent PDE. The diffusion term is assumed to be moderately stiff, while the reaction term can be very stiff. The advection can be such that the Peclet number is $\bigo(1)$. If we proceed in the same way as for \eqref{eq:ode}, we get ($a_{ii}=\gamma, i=1,\ldots,m$)
%\begin{equation}\label{eq:adrprox}
%\begin{split}
%\dot{\boldsymbol{Y}} &=
 % \Delta t \gamma F_D[\boldsymbol{Y}] + (\boldsymbol{Y}^0 - \boldsymbol{Y}) \\
%  &+ \Delta t \big((\calA - \gamma \mathbbm{1}_m) \otimes \mathbbm{1}_d\big) F_D[\boldsymbol{Y}] + \Delta t (\calA \otimes \mathbbm{1}_d) F_A[\boldsymbol{Y}] \\ 
%  &+ \Delta t (\calA \otimes \mathbbm{1}_d) F_R[\boldsymbol{Y}] \\
%  &\coloneqq G_D(\boldsymbol{Y}) + G_A(\boldsymbol{Y}) + G_R(\boldsymbol{Y}).
%\end{split}
%\end{equation}
 
 %Consider the following recurrence (preserves steady state) to iterate on \eqref{eq:adrprox}
	
%	\begin{equation}
%		\begin{split}
%			G_0&=J^{-1}(G_D(\boldsymbol{Y}_0)+G_A(\boldsymbol{Y}_0)+G_R(\boldsymbol{Y}_0)),\\
%			K_0&=\boldsymbol{Y}_0,\\
%			K_1&=\boldsymbol{Y}_0+\mu_1hG_0,\\
%			K_j&=\mu_jh(G_D(K_{j-1})+G_0-G_D(\boldsymbol{Y}_0))+\nu_jK_{j-1}+(1-\nu_j)K_{j-2},\quad j=2,\dots,s,\\
%			\boldsymbol{Y}_1&=K_s,
%		\end{split}
%	\end{equation}
%	where $J=I_{md}-h G_R'(\boldsymbol{Y}_0)$, $h$ is the optimization step size, and the coefficients are those of the Chebyshev method. 

%    For linear problems, we obtain the following amplification factor (see the proofs of Theorems \ref{thm:diff} and \ref{thm:ad})
 %   \begin{equation*}
  %      M= \mathbbm{1}_m \otimes R(z_1) + (\calA-\gamma \mathbbm{1}_m) \otimes B(z_1)z_2+\calA\otimes B(z_1)q +\calA \otimes B(z_1)r,
   % \end{equation*}
   % where $r=h\Delta t \Xi (\mathbbm{1}_d-h\Delta t \Xi)^{-1}$, and $\Xi$ represents the linear reaction matrix.
    %For linear diffusion, the above method gives the following relation after one step
	%\begin{equation}\label{eq:stabfunadr}
		%\boldsymbol{Y}_1=\boldsymbol{Y}_0+\frac{\frac{T_s(\omega_0+\omega_1hD)}{T_s(\omega_0)}-1}{hD}hG_0.
	%\end{equation}

% \textcolor{red}{REWRITE THE STABILITY FUNCTION AS $A(z)+B(z)G_0$ AND EXPLAIN WHY THIS APPROACH WORKS. CONSIDER LINEAR DIFFUSION ONLY AND SHOW CONVERGENCE.}

% \textcolor{blue}{
% Modified scheme (added by Gilles on 09.10)
% For each internal stage index $i=1,\ldots,s$, consider the splitting (depending on $i$),
% 	\begin{align}
% 		\dot Y_i&=\left(h \gamma 
% 		F_D(Y_i)-Y_i+Y^0
%   \right)
% 	+
%  \left(h \sum_{j=1}^{i-1}
% a_{ij} F_D(Y_j)
% + h \sum_{j=1}^{i}
% a_{ij} F_A(Y_j)
%  \right)
% +h \sum_{j=1}^{i}
% a_{ij} F_R(Y_j)
% 		\label{eq:adrprox} \\
%  &:=G_{D,i}(\boldsymbol{Y})+G_{A,i}(\boldsymbol{Y})+G_{R,i}(\boldsymbol{Y}). \nonumber
% 	\end{align}
% or the variant (for a cheaper Jacobian w.r.t. the reaction part):
% 	\begin{align}
% 		\dot Y_i&=\left(h \gamma 
% 		F_D(Y_i)-Y_i+Y^0
%   \right)
% 	+
%  \left(h \sum_{j=1}^{i-1}
% a_{ij} (F_D(Y_j)+F_R(Y_j))
% + h \sum_{i=1}^{j}
% a_{ij} F_A(Y_j)
%  \right)
% +h \gamma F_R(Y_i)
% 		\label{eq:adrprox} \\
%  &:=G_{D,i}(\boldsymbol{Y})+G_{A,i}(\boldsymbol{Y})+G_{R,i}(\boldsymbol{Y}). \nonumber
% 	\end{align}
%  }

  
%\subsection{Some idealised analysis}

  
	\section{Numerical tests}
 \label{sec:num}
In this section, we present a number of  numerical experiments on the Brusselator problem for the SDIRK4 method. More precisely, we study the case of a highly stiff reaction as well as a large advection setting.  We denote with exSDIRK4  the proposed explicit stabilised implementation of SDIRK4, with \textsc{Matlab} prototype codes made publicly available at \cite{codesMatlabexSDIRK26}, and we compare its performance with the order 2 PIROCK method. Note that established explicit methods of order $4$ such as ROCK4 \cite{Abd02}, are not suitable for the two problems studied here due to the stiff nature of the advection and reaction terms. In all our numerical experiments, the time step $\Delta t$ is chosen adaptively as described in the remark below. 
\begin{remark} 
Compared to the standard embedded error estimator for stiff problems recalled in Remark \ref{rem:variable_step}, see also \cite[Chap. IV.6]{HaW96}, the only difference for the explicit stabilised implementation is that the inverse Jacobian term $J^{-1}$ involved in \eqref{eq:shampine} is no longer available. This is indeed the main feature of the proposed approach to avoid such a possibly large dimension Jacobian and its inverse, which can raise linear algebra issues. Alternatively to \eqref{eq:shampine}, we implement the Shampine trick using the following error estimator of the exSDIRK4 method,
$$
err={err}_s
$$
where we define $\overline {err}= J^{-1}(y_1-\hat y_1)$ where $J = I_{d} - h\, \Delta t \gamma F'_R(y_1)$ is the Jacobian involved for the reaction term, and the ${err}_j$ are computed by induction on $j\in \{1,\cdots,s\}$ using ${err}_0 = {err}_{-1} = 0$ as
$${err}_j = \mu_j \left(h \Delta t \gamma (F_D(y_1+{err}_{j-1}) - F_D(y_1))+ \overline{err}\right)+
\nu_j {err}_{j-1}-(\nu_j-1) {err}_{j-2},$$ 
where $\mu_1 = \omega_1/\omega_0, \nu_1=1$ and $\mu_j,\nu_j$ and the stage number $s$ are defined in Algorithm \ref{alg:3} (see \eqref{eq:rkcd_iter}).
We emphasize that the above error estimator benefits from $L$-stability with respect to diffusion and reaction terms in the spirit of the the standard estimator \eqref{eq:shampine}. Indeed, firstly in the absence of diffusion and reaction ($F_D=F_A=0$) it reduces to $err=\overline{err}=J^{-1}(y_1-\hat y_1)$ which exactly coincides with \eqref{eq:shampine}, and secondly in the case of a linear diffusion $F_D=Dy$, it yields $err_s=B_s(h\Delta t \gamma D) J^{-1} (y_1-\hat y_1)$ analogously to \eqref{eq:it_rkcd1} and where we recall the bound $|B_s(z)|\leq C/|z|$.
\end{remark}


%We will perform numerical experiments on the Brusselator problem considered in the PIROCK paper \cite{AV13}. The time step $\Delta t$ is chosen \emph{adaptively}
%using the embedded error estimator for SDIRK4 described in \cite[Chap. IV.6]{HaW96}.

 
Throughout the numerical experiments we make the following implementation assumption: the evaluations of $F_D[\boldsymbol Y]$, $F_A[\boldsymbol Y]$, and $F_R[\boldsymbol Y]$ for the $m$ internal
stages, are computed in \emph{parallel}. This is natural given the
block-diagonal structure of the stabilisation operators $R_s(h\mathcal{V}_1)$ and
$B_s(h\mathcal{V}_1)$, which decouple across stages, and reflects the remark following
Theorem~\ref{thm:diff}. For a vector $v\in R^d$ that represents a discretized function $f:\R^d\rightarrow\R^d$, we define the discrete $L^2-$norm by
\[\|v\|_{L^2} = \sqrt{\frac{\sum_{i=1}^{d}v_i^2}d}.\]
\begin{remark}
    As we mentioned before, for simplicity, we use $\mathbbm{1} \otimes y_n$ as initial guess in Algorithm \ref{alg:3} and, to coincide with the true solution of SDIRK4, we choose a fixed $tol_{it} = 10^{-12}$ as optimisation tolerance. However, the cost of exSDIRK4 discussed in the numerical tests below can be further reduced by choosing a better initial guess using interpolation or dense output techniques, and by choosing $tol_{it}$ smaller than the relative tolerance by only one or two orders of magnitude.
\end{remark}
	\subsection{1D Brusselator problem with highly stiff reaction}
    
	\begin{table}[b!]
		\centering
		\begin{tabular}{|c|c|c|c|c|c|c|}
			\hline
			Method & $tol$   &  $F_D$ evals & $F_R$ evals &  $F'_R$ evals &  acc(rej) & error \\
			\hline
			exSDIRK4 & $10^{-1}$ & $4447$ & $245$ & $40$ & $20(0)$ & $3.5\times 10^{-4}$\\
			PIROCK & $10^{-1}$ & $548$ & $63$ & $12$ & $12(0)$ & $2\times 10^{-2}$ \\
			\hline
			exSDIRK4 & $10^{-3}$ & $5072$ & $559$ & $77$ & $42(3)$ &$2.0\times 10^{-5}$ \\
			PIROCK & $10^{-3}$ & $870$ & $123$ & $24$ & $24(0)$ & $3.3\times 10^{-4}$ \\
			\hline
			exSDIRK4 & $10^{-5}$ & $7718$ & $1307$ & $183$ & $102(3)$ & $1.4\times10^{-6}$ \\
			PIROCK & $10^{-5}$ & $5719$ & $6205$ & $1239$ & $1234(6)$ & $9.9\times10^{-6}$\\
			\hline
			exSDIRK4 & $10^{-7}$ & $13463$ & $3053$ & $439$ & $231(4)$ & $1.7\times10^{-8}$\\
			PIROCK & $10^{-7}$ & $41608$ & $54967$ & $10992$ & $10987(6)$ & $1.4\times10^{-7}$\\
			\hline
			exSDIRK4 & $10^{-9}$ & $23056$ & $6248$ & $915$ & $468(5)$ & $2.2\times10^{-10}$\\
			PIROCK & $10^{-9}$ & $191700$ & $279447$ & $55888$ & $55888(2928)$ & $1.0\times10^{-9}$\\
			\hline
		\end{tabular}
		\captionof{table}{{Comparison between our explicit stabilised implementation of SDIRK4 and PIROCK for different} {prescribed tolerances.} }
		\label{table:bruss}
	\end{table}
	Consider the following Brusselator problem
	\begin{equation}\label{eq:bruss1d}
		\begin{split}
		\partial_t u &= \alpha\Delta u + A +u^2v - (B+1)u,\\
		\partial_t v &= \alpha\Delta v -u^2v + Bu,\\
        u(0,t)&=u(1,t), \quad v(0,t)=v(1,t),
		\end{split}
	\end{equation}
	where $x\in (0,1)$, $t\in (0,T)$, $\alpha=0.2$, $A=1$, $B=3\times 10^7$, and $T=1$. We consider the following initial condition
	$$
	u(x,0)=1+\sin(2\pi x) \quad \text{and} \quad v(x,0)=3, 
	$$
    and periodic boundary conditions.
	We use second order finite differences for the space discretisation with mesh size $\Delta x =1/200$, and the starting time step $\Delta t_0=10^{-6}$. %In Table \ref{table:bruss}, we compare our explicit stabilised implementation of SDIRK4 \eqref{eq:sdirk4} with PIROCK \cite{AV13}, the state-of-the-art method for stiff advection-diffusion-reaction equations.


    \begin{figure}[htbp]
    \centering
    % First subfigure
    \begin{subfigure}[b]{0.325\textwidth}
        \centering
        \includegraphics[width=\textwidth]{bruss1d_fd.png}
        \caption{$F_D$ evaluations.}
        % \label{fig:bruss1d_fd}
    \end{subfigure}
    \hfill
    % Second subfigure
    \begin{subfigure}[b]{0.325\textwidth}
        \centering
        \includegraphics[width=\textwidth]{bruss1d_fr.png}
        \caption{$F_R$ evaluations.}
        % \label{fig:bruss1d_fr}
    \end{subfigure}
    \hfill
    % Third subfigure
    \begin{subfigure}[b]{0.325\textwidth}
        \centering
        \includegraphics[width=\textwidth]{bruss1d_jac.png}
        \caption{{$F_R'$ evaluations}}
        % \label{fig:bruss1d_frjac}
    \end{subfigure}

    \caption{Cost comparison between exSDIRK4 and PIROCK for Problem \eqref{eq:bruss1d}}
    \label{fig:bruss1d}
\end{figure}

Table~\ref{table:bruss} and Figure~\ref{fig:bruss1d} report the cost comparison between exSDIRK4 and PIROCK for the 1D Brusselator problem {\eqref{eq:bruss1d}}. For loose tolerances ($tol = 10^{-1}$ and $10^{-3}$), PIROCK requires fewer function evaluations overall; however, exSDIRK4 already delivers a substantially lower error for the same requested tolerance, reflecting its higher order of accuracy. As the tolerance tightens, the balance shifts decisively in favour of exSDIRK4. At $tol = 10^{-5}$, both methods have a comparable number of $F_D$ evaluations, but exSDIRK4 requires more than $4$ times fewer $F_R$ evaluations and $6$ times fewer $F'_R$ {evaluations/inversions}. This advantage grows rapidly: at $tol = 10^{-7}$, PIROCK needs approximately $3$ times more $F_D$, $18$ times more $F_R$, and $24$ times more $F'_R$ {evaluations/inversions} than exSDIRK4. At the tightest tolerance $tol = 10^{-9}$, PIROCK accumulates 55,888 accepted steps with 2,928 rejections, while exSDIRK4 completes the integration in only 468 accepted steps with 5 rejections—a reduction of more than two orders of magnitude in the number of steps. These trends are clearly visible in Figure~\ref{fig:bruss1d}, which displays the number of $F_D$, $F_R$, and $F_R'$ evaluations as a function of the tolerance. The work-precision curves show that the crossover point beyond which exSDIRK4 is uniformly cheaper occurs around $tol \approx 10^{-4}$, and the gap widens as the tolerance decreases, which is consistent with the high-order convergence of exSDIRK4.


\subsection{2D Brusselator with large advection}
{We now consider the following Brusselator model in 2D:}
\begin{equation}\label{eq:bruss2d}
		\begin{split}
		\partial_t u &= \nu\Delta u + \mu U\cdot \nabla u + A +u^2v - (B+1)u,\\
		\partial_t v &= \nu\Delta v + \mu V\cdot \nabla v -u^2v + Bu,\\
		\end{split}
	\end{equation}
	where $x\in (0,1)^2$, $t\in (0,T)$, $\nu=0.1$, $\mu=2$ $A=1.3$, $B=10^7$, $U=(-0.5,1)^T$, $V=(0.4,0.7)^T$ and $T=0.5$. Following \cite{AV13}, we choose
	$
	u(x,0)=22x_2(1-x_2)^{3/2} , \quad v(x,0)=27x_1(1-x_1)^{3/2}, 
	$
	and periodic boundary conditions. For the space discretisation, we use again second order finite differences with mesh size equals to $10^{-2}$ in both dimensions and the starting time step $\Delta t_0=10^{-6}$.

\begin{figure}[htbp]
    \centering
    % First subfigure
    \begin{subfigure}[b]{0.325\textwidth}
        \centering
        \includegraphics[width=\textwidth]{bruss2d_fd.png}
        \caption{$F_D$ evaluations.}
        % \label{fig:bruss2d_fd}
    \end{subfigure}
    \hfill
    % Second subfigure
    \begin{subfigure}[b]{0.325\textwidth}
        \centering
        \includegraphics[width=\textwidth]{bruss2d_far.png}
        \caption{$F_A$ and $F_R$ evaluations.}
        % \label{fig:bruss2d_far}
    \end{subfigure}
    \hfill
    % Third subfigure
    \begin{subfigure}[b]{0.325\textwidth}
        \centering
        \includegraphics[width=\textwidth]{bruss2d_jac.png}
        \caption{{$F_R'$ evaluations.}}
        % \label{fig:bruss2d_frjac}
    \end{subfigure}
    \caption{Cost comparison between exSDIRK4 and PIROCK for Problem \eqref{eq:bruss2d}}
    \label{fig:bruss2d}
\end{figure}

Figure~\ref{fig:bruss2d} reports analogous results for the 2D Brusselator problem with large advection. The qualitative behaviour is very similar to the 1D case: exSDIRK4 is outperformed by PIROCK for loose tolerances in terms of raw function evaluation counts, but becomes significantly more efficient as the tolerance tightens. The three panels—$F_D$ evaluations, $F_A$/$F_R$ evaluations, and $F_R'$ evaluations—all exhibit the same crossover pattern observed in Figure~\ref{fig:bruss1d}, confirming that the efficiency advantage of exSDIRK4 is not an artefact of the one-dimensional setting but carries over to higher-dimensional problems with combined diffusion, advection, and stiff reaction terms.

\section{Conclusion}

%\paragraph{Conclusion}

We introduced an explicit stabilised implementation of singly diagonally implicit
Runge-Kutta (SDIRK) methods for advection-diffusion-reaction PDEs, with the fourth-order
method SDIRK4 as the primary example. The central idea is to reformulate the implicit
Runge-Kutta stage equations \eqref{eq:nonlinearimp} as the steady state of an auxiliary
ODE \eqref{eq:gf1}, and to compute that steady state by a partitioned
Runge-Kutta-Chebyshev iteration (Algorithm~\ref{alg:3}) inspired by the RKCD
optimisation method of \cite{EVVZ21}. The partition of the Runge-Kutta coefficient matrix
$\calA = \gamma I_m + (\calA - \gamma I_m)$ is essential: it isolates a symmetric,
block-diagonal stiff part $G_D$ that governs the stability of the iteration through the
stability polynomial $R_s$, avoiding the non-diagonalisable amplification matrix that
would otherwise prevent convergence, as analysed in detail in Section~\ref{appendix:naive}.
The off-diagonal diffusion and advection contributions enter through the damping polynomial
$B_s$, whose extended stability along the imaginary axis accommodates large advection
velocities under a mild CFL-like condition on the iteration step size $h$ — significantly
more permissive than the classical CFL restriction on the physical time step $\Delta t$.
The stiff reaction term is handled implicitly at each inner iteration via a local inversion
of $J = I_{md} - hG'_R(\boldsymbol{x}_k)$, which remains computationally inexpensive
because the reaction Jacobian carries no spatial coupling between grid points and is
therefore block-diagonal. The result is a method that requires no large-scale Newton solves
or global linear algebra, and yet achieves the high-order accuracy and strong stability of
a fourth-order $L$-stable integrator. The only implicit component is this local reaction
Jacobian inversion, in sharp contrast to the dense, high-dimensional Jacobian systems that
a standard Newton-based implementation of SDIRK4 would require.

The convergence and stability properties of the new method are rigorously established in
two settings. In the pure linear diffusion case (Theorem~\ref{thm:diff}), we prove that
Algorithm~\ref{alg:3} converges with {decay factor} $\alpha_s(\eta)+\varepsilon$ for any
$\varepsilon>0$, independently of the spatial dimension $d$, and that implementing one
step of an $m$-stage SDIRK method requires only
$\mathcal{O}(\sqrt{\Delta t}/\Delta x)$ evaluations of $F_D$ — the same asymptotic cost
as a standard first-order explicit stabilised method. In the linear advection-diffusion
case (Theorem~\ref{thm:ad} and Propositions~\ref{prop:rate1}--\ref{prop:rate2}), we show
that convergence is maintained provided the advection-to-diffusion ratio satisfies a
mild CFL-like condition, a condition that is significantly
more permissive than the one imposed by a classical explicit scheme on the physical time
step $\Delta t$.

Assuming parallel implementation of the function evaluations across the method stages,
the numerical experiments on the 1D and 2D Brusselator problems with {highly stiff
reaction  confirms} the theoretical predictions and demonstrate a clear
advantage over the second-order method PIROCK \cite{AV13} whenever tight tolerances are
required. The method occupies a natural middle ground: more implicit than a pure explicit
stabilised method such as PIROCK, but far cheaper than a fully Newton-based implicit
solver, and this balance is precisely what drives its efficiency at high accuracy demands.

%\paragraph{Future Work}
%Several directions for future research naturally follow from the present work.
%\paragraph{Extension to fully implicit Runge-Kutta methods.}
The present framework relies on the lower-triangular, singly diagonal structure of SDIRK
methods to decouple the implicit system stage by stage and to identify the symmetric
stabilising part $G_D$. A natural and challenging next step would be to extend the approach to
fully implicit Runge-Kutta methods such as Radau IIA or Gauss-Legendre collocation
schemes. 
%These methods achieve the theoretical maximum order for $m$ stages and are known for their superior damping properties. 
{A major difficulty for their implementation is their full matrix structure of the Runge-Kutta coefficients $\calA$, and studied for instance in \cite{IM98,BIM15} where splitting approaches based on  triangular matrices and low rank linear algebra techniques are proposed. 
In such fully implicit case,} the
Runge-Kutta matrix $\calA$ is a full, non-triangular matrix with non-real eigenvalues,
so the partitioning strategy proposed in this paper and the associated convergence analysis do not apply straightforwardly.

\paragraph{Acknowledgments} 
The authors would like to thank Ernst Hairer and Francesca Mazzia for helpful discussions on the implementation of SDIRK and RADAU methods. G.V. was partially supported by the Swiss National Science
Foundation, projects No. 200020\_214819, No. 200020\_192129 and No. 10009199. I.A. was partially supported by the Swiss National Science Foundation, projects No. P500PT\_210968 and No. P5R5PT\_222213.


   \bibliographystyle{abbrv}
	\bibliography{refs}

%\paragraph{Nonlinear convergence analysis.}
%The convergence results of Section~4 are established for linear problems. While the
%numerical experiments suggest that Algorithm~\ref{alg:3} performs well on nonlinear
%problems such as the Brusselator, a rigorous nonlinear convergence theory is currently
%missing. Extending the analysis to mildly nonlinear or Lipschitz-continuous vector fields,
%potentially via contraction mapping arguments or Lyapunov-based techniques, would
%strengthen the theoretical foundations of the method and provide sharper parameter
%selection guidelines.

%\paragraph{Applications to real-world problems.}
%The method is directly applicable to a broad class of high-dimensional stiff PDEs
%that arise in science and engineering. Relevant application areas include:
%\begin{itemize}
%    \item \emph{Plasma physics:} Transport equations for ions and electrons in fusion
%    devices (e.g.\ tokamaks) involve multi-scale diffusion-advection-reaction dynamics
%    in up to three spatial dimensions, with highly anisotropic diffusion coefficients and
%    stiff source terms. The high-order accuracy and dimension-independent cost of
%    exSDIRK4 make it a natural candidate for such problems.
%    \item \emph{Combustion and reacting flows:} Large chemical mechanisms in reactive
 %   flow simulations give rise to extremely stiff reaction terms coupled to convection and
%    diffusion. The ability of exSDIRK4 to handle stiff reactions via a cheap block-diagonal
%    Jacobian inversion, without requiring large-scale Newton solves, is particularly
 %   attractive in this setting.
%    \item \emph{Optimal control and PDE-constrained optimisation:} The connection between
  %  explicit stabilised methods and optimisation algorithms exploited in \cite{EVVZ21} and
 %   in the present work suggests further synergies, for instance in the time integration of
 %   adjoint equations arising in optimal control of parabolic PDEs \cite{AV21}.
%\end{itemize}
\appendix

\section{Naive implementation}
\label{appendix:naive}
We now explain in more detail the issues observed in Figure \ref{fig:iterations} with respect to the naive stabilised explicit implementation of SDIRK4 applied without partitioning {(corresponding to Algorithm \ref{alg:3} with 
the definition of $G(\boldsymbol{x})$ replaced by $G(\boldsymbol{x})=G_D(\boldsymbol{x})+G_A(\boldsymbol{x})$ with $G_R(\boldsymbol{x})=0$, or equivalently applying Algorithm \ref{alg:rkcd} to equation \eqref{eq:gff} with $F_A=F_R=0$)}. 
Considering for simplicity only one mode of the discrete linear heat equation, corresponding to the Dahlquist test equation $\dot y =\lambda y$ 
with $\lambda=-\Delta t \gamma \lambda_{\text{heat},j}<0$ corresponds to a possibly large parameter $|\lambda|\gg 1$.
We then write one-step recursion induced by Algorithm \ref{alg:rkc} applied to a SDIRK method with $m>1$ internal stages, namely  %For simplicity, we will restrict ourselves to the case of $N=1$, but we insist that the analysis applies to higher dimensional problems. In this case, using Algorithm \ref{alg:rkc} we have the following recursion relation
\[
\mathbf{Y}_{n+1}=R_{s}(-hI_m+ \lambda h\mathcal{A})\mathbf{Y}_{n}
\]
where $R_{s}$  given by \eqref{eq:stabcheb} is the stability function of RKC, and $h$ is the time step used by the RKC iteration. 
%Recall that we search for a bound on $\norm{R_{s}(Qh)}$ but because RKC is not an A-stable method and $Q$ in the case of SDIRK4 is not diagonalisable implies using the fact that the eigenvalues of $Q$ are negative is not enough. 
Consider the Jordan form of the matrix $\mathcal{A}=P\mathcal{J}P^{-1}$ where
$\mathcal{J}$ is a Jordan block defined in \eqref{eq:defJ}.
%        $$
%    \mathcal{A} = \begin{pmatrix}
%	\gamma\\
%	& \gamma\\
%	& (a_{ij})_{i>j} &  \ddots\\
%	& & & \gamma \\
%    \end{pmatrix}, \qquad
%    \calJ=\begin{pmatrix}
%	\gamma&1\\
%	& \gamma&1\\
%	  & & \ddots&\ddots\\
%	& & & \gamma& 1 \\
%        &&&&\gamma
%    \end{pmatrix}.
%        $$
Using the following formula for any analytic function $f$ applied to the Jordan matrix $J$,
        $$
f(\calJ)=\begin{pmatrix}f(\gamma)&f^\prime(\gamma)&\cdots&\frac{f^{(m-1)}(\gamma)}{(m-1)!}\\&f(\gamma)&\ddots&\vdots\\&&\ddots&f^\prime(\gamma)\\0&&&f(\gamma)\end{pmatrix},
        $$
        we deduce taking $f(\calJ)=R(-hI_m+\gamma z \calJ)$ with $z=h\lambda$ and $\hat z=-h+\gamma z$,
$$
P^{-1}R_s(-hI_m +\lambda h \mathcal{A})P = \begin{pmatrix}
R_s(\hat z)&R_s^\prime(\hat z) \gamma z&\cdots& \frac{R_s^{(m-1)}(\hat z)(\gamma z)^{m-1}}{(m-1)!} \\&R_s(\hat z)&\ddots&\vdots\\&&\ddots&R_s^\prime(\hat z) \gamma z\\0&&&R_s(\hat z)
\end{pmatrix}
$$
and it turns out that the off-diagonal terms $\frac{R_s^{(j)}(\hat z)(\gamma z)^{j}}{j!}, 1\leq j\leq m-1$ are arbitrarily large for large $|z|$ for stiff problems where $|\lambda| \gg 1$. 
This is due to the $z^j$ factors, growing like $\bigo(s^{2j})$, and can be shown that the 
derivatives $R_s^{(j)}$ of the stability function oscillate with an amplitude of size $\bigo(s^{-j})$ which is not sufficient to damp the $z^j$ terms.
We obtain  
$$\sup_{s\geq 1, \hat z\in [-L_{s,\eta},0]} \|R_s(-hI_m+z \mathcal A)\|=+\infty$$
and this makes the convergence factor of the naive implementation without partitioning non-uniformly bounded 
as the stiffness increases when refining the spatial mesh ($\Delta x\rightarrow 0$), as illustrate in Figure \ref{fig:iterations}.
Note that in the proof of the main Theorem \ref{thm:diff}, we used the same Jordan decomposition of the Runge-Kutta matrix
$
\calA = P\mathcal{J}P^{-1}
$, with the main difference in the design of the proposed method with partitioning and its stability analysis that the stability function $R_s$ of the explicit stabilised Chebyshev method is always applied to a diagonalizable matrix.

	\end{document}
